\exercisegroup

\begin{enumerate}
\def\labelenumi{\arabic{enumi})}
\tightlist
\item
  设 E 与 F 是两个满足第一可数公理的无穷维希尔伯特空间，\((a_{n})\) 是 E 的一个标准正交基，\((b_{n})\) 是 F 的一个标准正交基。
\end{enumerate}

\begin{enumerate}
\def\labelenumi{\alph{enumi})}
\tightlist
\item
  设 u 是从 E 到 F 中的连续线性映射；设 \(u(a_{n}) = \sum_{m} \alpha_{mn} b_{m}\)。证明
\[
\sum_ {n} | \alpha_ {m n} | ^ {2} \leqslant \| u \| ^ {2}
\]

\[
\sum_ {m} | \alpha_ {m n} | ^ {2} \leqslant \| u \| ^ {2}
\]
对每个 \(m\) 与 \(n\) 成立。

\begin{enumerate}
\def\labelenumi{\alph{enumi})}
\setcounter{enumi}{1}
\tightlist
\item
  给出一个二重序列 \((\alpha_{mn})\) 的例子，使得 \(\sum_{m} |\alpha_{mn}|^{2} \leqslant 1\) 对所有 \(n\) 成立，且 \(\sum_{n} |\alpha_{mn}|^{2} \leqslant 1\) 对所有 \(m\) 成立，但不存在任何从 E 到 F 中的连续线性映射 \(u\)，使得 \(\langle u(a_n)|b_m\rangle = \alpha_{mn}\) 对每对整数 \((m, n)\) 成立。（证明：若 I \(\subset\) N 是由 \(p\) 个整数组成的集，且 V \(_p\)（相应地 W \(_p\)）是 E（相应地 F）中由 \(a_n\)（相应地 \(b_n\)，均满足 \(n \in I\)）生成的子空间，则存在一个线性映射 \(u_p\)（从 V \(_p\) 到 W \(_p\) 上的），使得 \(\langle u_p(a_n)|b_m\rangle = \frac{1}{\sqrt{p}}\) 对 \(m \in I\) 与 \(n \in I\) 成立，且 \(\|u_p\| \geqslant \sqrt{p}\)。）
\end{enumerate}

\begin{enumerate}
\def\labelenumi{\arabic{enumi})}
\setcounter{enumi}{1}
\tightlist
\item
  设 \(A = (\alpha_{mn})_{(m,n)\in \mathbf{N}\times \mathbf{N}}\) 是一个复数的二重序列，我们也称它为无穷矩阵。对每个点 \(x = (x_{n})\)（直和空间 \(\mathbf{C}^{(\mathbf{N})}\) 中的），和式 \(y_{m} = \sum_{n}\alpha_{mn}x_{n}\) 都有定义；令 \(A.x\) 表示点 \((y_{m})\)，它属于乘积空间 \(\mathbf{C}^{\mathbf{N}}\)，于是 \(x\mapsto A.x\) 是从 \(\mathbf{C}^{(\mathbf{N})}\) 到 \(\mathbf{C}^{\mathbf{N}}\) 中的一个线性映射，且从 \(\mathbf{C}^{(\mathbf{N})}\) 到 \(\mathbf{C}^{\mathbf{N}}\) 中的每个线性映射都具有这一形式。用 \(E_{n}\) 表示 \(\mathbf{C}^{(\mathbf{N})}\) 中由典范基前 \(n\) 个向量生成的子空间，用 \(P_{n}\) 表示从 \(\mathbf{C}^{\mathbf{N}}\) 到 \(E_{n}\) 上的典范投影；当赋给 \(E_{n}\) 由空间 \(\ell_{\mathbf{C}}^{2}(\mathbf{N})\) 的范数诱导的范数时，\(\| u\|\) 表示线性映射 \(u\)（从有限维希尔伯特空间 \(E_{n}\) 到其自身的）的范数。
\end{enumerate}

\begin{enumerate}
\def\labelenumi{\alph{enumi})}
\item
  为使 \(\mathbf{C}^{(\mathbf{N})}\) 在映射 \(x\mapsto A.x\) 下的像含于 \(\ell_{\mathbf{C}}^2 (\mathbf{N})\)，且该映射可延拓为从 \(\ell_{\mathbf{C}}^2 (\mathbf{N})\) 到其自身的连续线性映射，必要且充分的是诸范数 \(\| P_nAP_n\|\) 有界。这蕴含 \(A\) 的行与列都属于 \(\ell_{\mathbf{C}}^2 (\mathbf{N})\)（习题 1）。
\item
  用 \(A^*\) 表示无穷矩阵 \((\alpha_{mn}')\)，其中 \(\alpha_{mn}' = \overline{\alpha}_{nm}\)。若 \(A\) 的列都属于 \(\ell_{\mathbf{C}}^2(\mathbf{N})\)（换言之，若 \(x \mapsto A.x\) 把 \(\mathbf{C}^{(\mathbf{N})}\) 映到 \(\ell_{\mathbf{C}}^2(\mathbf{N})\) 中），则级数 \(\beta_{mn} = \sum_{p} \overline{\alpha}_{pm} \alpha_{pn}\)
\end{enumerate}

绝对收敛，此时我们置 \(A^{*}A = (\beta_{mn})\)。证明：为使 \(x \mapsto A.x\) 可延拓为一个连续线性映射 \(u\)（从 \(\ell_{\mathbf{C}}^{2}(\mathbf{N})\) 到其自身的），必要且充分的是诸范数 \(\| P_n(A^*A)P_n\|\) 有界（我们有 \(\langle P_n(A^*A)P_n.x|x\rangle = \| AP_n .x\|^2\)，对所有 \(x\in E_n\) 成立）。于是 \(x\mapsto A^{*}A.x\) 延拓为一个正埃尔米特映射 \(u^{*}u\)（从 \(\ell_{\mathbf{C}}^{2}(\mathbf{N})\) 到其自身的）。

\begin{enumerate}
\def\labelenumi{\alph{enumi})}
\setcounter{enumi}{2}
\item
  对两个无穷矩阵 \(X = (\xi_{mn})\)、\(Y = (\eta_{mn})\)，我们称乘积 \(XY\) 有定义，如果级数 \(\zeta_{mn} = \sum_{p} \xi_{mp} \eta_{pn}\) 绝对收敛，此时我们置 \(XY = (\zeta_{mn})\)。我们称幂 \(X^k\)（\(k\) 为整数 \(> 1\)）有定义，如果 \(X^{k-1}\) 与 \(X^{k-1}X\) 都有定义，此时我们置 \(X^k = X^{k-1}X\)；在这一情形，\(X^pX^q = X^k\) 对每对整数 \(p, q\)（满足 \(p + q = k\)）成立。设 \(A\) 是一个列都属于 \(\ell_{\mathbf{C}}^2(\mathbf{N})\) 的无穷矩阵，且乘积 \((A^*A)^2\) 有定义，证明：对所有 \(x \in E_n\)，我们有 \(\langle (P_nA^*AP_n)^2.x|x\rangle \leqslant \langle P_n(A^*A)^2P_n.x|x\rangle\)，并由此推出 \(\| P_nA^*AP_n\|^2 \leqslant \| P_n(A^*A)^2P_n\|\)。
\item
  为使无穷矩阵 A 满足：\(\mathbf{C}^{(\mathbf{N})}\) 在 \(x\mapsto A.x\) 下的像含于 \(\ell_{\mathbf{C}}^2 (\mathbf{N})\)，且 \(x\mapsto A.x\) 可延拓为从 \(\ell_{\mathbf{C}}^2 (\mathbf{N})\) 到其自身的连续线性映射，必要且充分的是下列三个条件都被满足：
\end{enumerate}

\begin{enumerate}
\def\labelenumi{(\roman{enumi})}
\item
  \(A\) 的行与列都属于 \(\ell_{\mathbf{C}}^{2}(\mathbf{N})\)；
\item
  幂 \((A^{*}A)^{k}\) 对每个整数 \(k > 1\) 都有定义；
\item
  我们有
\end{enumerate}

\[
\sup _ {n} \left(\sup _ {m} | ((A ^ {*} A) _ {m m} ^ {n}) ^ {1 / n} |\right) <   + \infty
\]

其中 \((A^{*}A)_{mm}^{n}\) 表示以 \((m,m)\) 为指标的项，它属于矩阵 \((A^{*}A)^{n}\)。（注意，若 \(C\) 是 \(\mathrm{E}_n\) 的一个正自伴自同态关于 \(\mathrm{E}_n\) 的典范基的矩阵，则 \(\| C\| \leqslant n\sup_{1\leqslant i\leqslant n}|C_{ii}|\)（考察 \(C\) 的迹并将 \(C\) 对角化）。利用在 \(c\) ) 中证明的不等式，证明

\[
\left\| P _ {n} A ^ {*} A P _ {n} \right\| \leqslant n ^ {2 - k} \sup _ {1 \leqslant i \leqslant n} \left| ((A ^ {*} A) _ {i i} ^ {2 k}) \right| ^ {2 - k}
\]

对每个整数 \(k > 1\) 成立，只要条件 (i)、(ii) 与 (iii) 都被满足。）

¶ 3) 设 \((a_{ij})_{(i,j)\in\mathbb{I}\times\mathbb{I}}\) 是一个可数无穷的复数二重族。假设存在两个数 \(\beta > 0\)、\(\gamma > 0\)，以及一个由 > 0 的数组成的族 \((p_i)_{i\in\mathbb{I}}\)，满足关系式

\[
\sum_ {i} p _ {i} | a _ {i j} | \leqslant \beta p _ {j}, \quad \sum_ {j} | a _ {i j} | p _ {j} \leqslant \gamma p _ {i}\tag{*}
\]

对所有 \(i, j\)（属于 \(\mathbb{I}\)）成立。
a) 证明存在一个连续自同态 \(u\)（\(\ell_{\mathbf{C}}^2 (\mathrm{I})\) 上的），其范数 \(\leqslant (\beta \gamma)^{1 / 2}\)，使得对所有 \(x = (x_{i})_{i\in \mathbf{I}}\)（属于 \(\ell_{\mathbf{C}}^2 (\mathrm{I})\)），有 \(u(x) = y\)，其中 \(y = (y_{i})\) 由 \(y_{i} = \sum_{j}a_{ij}x_{j}\) 给出（对 \(x = (x_{i})\) 与 \(y = (y_{i})\)（\(\mathbf{C}^{(\mathbb{I})}\) 中的），置 \(v_{ij} = |x_i| (p_j|a_{ij}| / p_i)^{1 / 2}\)、\(w_{ij} = |y_j|(p_i|a_{ij}| / p_j)^{1 / 2}\)，并给出 \(\sum_{i,j}v_{ij}w_{ij}\) 的一个上界）。

\begin{enumerate}
\def\labelenumi{\alph{enumi})}
\setcounter{enumi}{1}
\tightlist
\item
  取 I 为所有整数 \(\geqslant 1\) 组成的集，并设 \(a_{ij} = (i + j)^{-1}\)。证明当 \(p_i = i^{-1/2}\) 且 \(\beta = \gamma = \pi\) 时条件 (*) 被满足（这些是可能取到的最佳常数）（把 (*) 中的级数与一个积分相比较）。以类似的方式处理 I = N 且 \(a_{ij} = (i + j + 1)^{-1}\) 的情形（「Hilbert 矩阵」）。
\end{enumerate}

* c) 设 \(\mathcal{H} = \mathrm{H}^2 (\mathbf{D})\) 是 Hardy 空间，它由所有满足下列条件的函数 \(f(z) = \sum_{n=0}^{\infty} a_n z^n\) 组成：在开圆盘 \(\mathbf{D}:|z| < 1\) 中解析，且 \(\| f\|^2 = \sum_n |a_n|^2 < +\infty\)；于是 \(\| f\|\) 是 \(\mathcal{H}\) 上的一个范数，对于它 \(\mathcal{H}\) 同构于 \(\ell_{\mathbf{C}}^2 (\mathbf{N})\)。给定两个函数 \(f, g\)（属于 \(\mathcal{H}\)），证明函数 \(t \mapsto f(t) g(t)\)（实变量 \(t\) 的）在 \([0,1]\) 上关于 Lebesgue 测度可积，且公式 \(\mathbf{B}(f,g) = \int_0^1 f(t) g(t) dt\) 在 \(\mathcal{H} \times \mathcal{H}\) 上定义一个连续双线性型（考察 \(\mathbf{B}(f,f) = \int_0^1 f(t)^2 dt\)，其中函数 \(f \in \mathcal{H}\) 形如 \(f(z) = \sum_{n=0}^{\mathbf{N}} a_n z^n\)，且 \(a_n \geqslant 0\)（对 \(0 \leqslant n \leqslant \mathbf{N}\)）；用 Cauchy 定理建立关系式

\[
\int_ {- 1} ^ {1} f (t) ^ {2} d t = - i \int_ {0} ^ {\pi} f (e ^ {i \theta}) ^ {2} e ^ {i \theta} d \theta
\]

由此得到 \(\mathbf{B}(f, f) \leqslant \frac{1}{2} \int_{-\pi}^{\pi} |f(e^{i\theta})|^2 d\theta = \pi \| f\|^2\)。用这一方法，得到 \(b\) ) 的结果：Hilbert 矩阵定义了一个范数 \(\leqslant \pi\) 的自同态（在希尔伯特空间 \(\ell_{\mathbf{C}}^2(\mathbf{N})\) 上）。*

\begin{enumerate}
\def\labelenumi{\arabic{enumi})}
\setcounter{enumi}{3}
\tightlist
\item
  设 \(\mathbf{E}\) 是一个维数为 \(d\) 的复希尔伯特空间。
\end{enumerate}

\begin{enumerate}
\def\labelenumi{\alph{enumi})}
\item
  对每个 \(u \geqslant 0\)（\(\mathcal{L}(\mathrm{E})\) 中的，V, p.~45），证明存在唯一的 \(v \geqslant 0\)（\(\mathcal{L}(\mathrm{E})\) 中的）使 \(u = v^2\)（把 \(u\) 对角化）；我们记 \(v = u^{1/2}\)。
\item
  对每个 \(u \in \mathcal{L}(\mathrm{E})\)，置 \(\mathrm{abs}(u) = (u^{*}u)^{1/2}\)。证明 u 与 \(\mathrm{abs}(u)\) 有相同的范数，且我们有 \(\mathrm{abs}(\symbf{\Lambda}^{n}(u)) = \symbf{\Lambda}^{n}(\mathrm{abs}(u))\)（对每个整数 \(n \leqslant d\)）。
\item
  设 \(s_{1}(u) \geqslant s_{2}(u) \geqslant \ldots \geqslant s_{d}(u) \geqslant 0\) 是 \(\mathrm{abs}(u)\) 的按重数计算的特征值序列。证明 \(\|u\| = s_{1}(u)\)，且对每个整数 \(n \leqslant d\)，\(\| \symbf{\Lambda}^{n}(u)\| = s_{1}(u) s_{2}(u) \ldots s_{n}(u)\)。为使 \(\| \symbf{\Lambda}^{n}(u)\| = \| u\|^{n}\)（对所有满足 \(1 \leqslant n \leqslant d\) 的 n），必要且充分的是 \(u^{*}u\) 是一个位似，换言之，u 是一个酉算子的标量倍。
\end{enumerate}

\begin{enumerate}
\def\labelenumi{\arabic{enumi})}
\setcounter{enumi}{4}
\tightlist
\item
  设 E、F 是两个希尔伯特空间，u 是从 E 到 F 中的连续线性映射。用 \(\ell(u)\) 表示所有 \(x \in E\)（满足 \(\|u(x)\| = \|u\|. \|x\|\)）组成的集。
\end{enumerate}

\begin{enumerate}
\def\labelenumi{\alph{enumi})}
\item
  证明 \(\ell(u)\) 是 E 的闭向量子空间，它是 \(u^{*}u - \| u\|^{2}1_{\mathrm{E}}\) 的核，且与 \(u\) 的核正交。
\item
  证明 \(u\) 在 \(\ell(u)\) 上的限制是从 \(\ell(u)\) 到 \(\ell(u^*)\) 上的双射，其逆双射是 \(\|u\|^{-2}.u^*\) 在 \(\ell(u^*)\) 上的限制；此外，正交补 \((\ell(u))^{\circ}\) 在 \(u\) 下的像含于 \((\ell(u^*))^{\circ}\)。若 \(u_1\) 是 \(u\) 在 \((\ell(u))^{\circ}\) 上的限制（视为从 \((\ell(u))^{\circ}\) 到 \((\ell(u^*))^{\circ}\) 中的映射），则其伴随 \(u_1^*\) 是 \(u^*\) 在 \((\ell(u^*))^{\circ}\) 上的限制；若 \(\ell(u) \neq E\)，用 \(\langle u\rangle\)（称为 \(u\) 的次范数）表示范数 \(\|u_1\|\)；若 \(\ell(u) = E\)，则置 \(\langle u\rangle = 0\)。于是 \(\langle u^*\rangle = \langle u\rangle\)。
\end{enumerate}

\begin{enumerate}
\def\labelenumi{\arabic{enumi})}
\setcounter{enumi}{5}
\tightlist
\item
  设 E 是希尔伯特空间，M、N 是 E 的两个闭子空间，\(M^{\circ}\)、\(N^{\circ}\) 是它们各自的正交补；用 \(p_{M}\)、\(p_{N}\) 分别表示到 M 与 N 上的正交投影，使 \(1_{E} - p_{M}\)、\(1_{E} - p_{N}\) 分别是到 \(M^{\circ}\) 与 \(N^{\circ}\) 上的正交投影。置 \(u_{\mathrm{NM}} = (1_{\mathrm{E}} - p_{\mathrm{N}}) p_{\mathrm{M}}\) 与 \(\delta(\mathbf{M}, \mathbf{N}) = \|u_{\mathrm{NM}}\|\)；我们有 \(\delta(\mathbf{M}, \mathbf{N}) = \delta(\mathbf{N}^{\circ}, \mathbf{M}^{\circ}) \leqslant 1\)；关系 \(\delta(\mathbf{M}, \mathbf{N}) < 1\) 蕴含 \(M \cap N^{\circ} = \{0\}\)。
\end{enumerate}

\begin{enumerate}
\def\labelenumi{\alph{enumi})}
\item
  用 \(\tilde{\mathbf{M}}\) 表示 \(\mathbf{M}\) 中 \(\mathbf{M} \cap \mathbf{N}^{\circ}\) 的正交补，并设 \(\varepsilon(\mathbf{M}, \mathbf{N}) = \delta(\tilde{\mathbf{M}}, \mathbf{N})\)。证明（用习题 5 的记号）\(\ell(u_{\mathrm{NM}}) = \mathbf{M} \cap \mathbf{N}^{\circ}\)，并由此推出 \(\varepsilon(\mathbf{M}, \mathbf{N}) = \langle u_{\mathrm{NM}} \rangle \leqslant \delta(\mathbf{M}, \mathbf{N})\)；此外，若 \(\mathbf{M} \cap \mathbf{N}^{\circ} = \{0\}\)（特别地，若 \(\delta(\mathbf{M}, \mathbf{N}) < 1\)），则 \(\varepsilon(\mathbf{M}, \mathbf{N}) = \delta(\mathbf{M}, \mathbf{N})\)。
\item
  对从 E 到其自身的连续线性映射 \(u\)，我们称 \(u\) 的余范数为数 \(c(u) = \inf \|u(x)\|/\|x\|\)，其中 \(x\) 取遍所有 \(\neq 0\) 且与 \(u^{-1}(0)\) 正交的向量组成的集（若 \(u = 0\)，置 \(c(u) = 1\)）。为使 \(u(\mathrm{E})\) 在 E 中是闭的，必要且充分的是 \(c(u) > 0\)（I, p.~17, th. 1）。我们有 \(c(u^{*}) = c(u)\)。
\item
  设 \(v_{\mathbf{NM}} = p_{\mathbf{N}}p_{\mathbf{M}}\)。证明
\end{enumerate}

\[
\varepsilon (\mathbf {M}, \mathbf {N}) ^ {2} + c (v _ {\mathbf {N M}}) ^ {2} = 1
\]

（注意 \(\|u_{\mathrm{NM}}(x)\|^{2} + \|v_{\mathrm{NM}}(x)\|^{2} = \|p_{\mathrm{M}} \cdot x\|^{2}\)，且 \(v_{NM}\) 的核是 \(M^{\circ} + (M \cap N^{\circ})\)，并由此推出 \(\langle u_{NM}\rangle^{2} \leqslant 1 - c(v_{\mathrm{NM}})^{2}\)。）
\item
  由 b) 与 c) 推出 \(\varepsilon(\mathrm{N}, \mathrm{M}) = \varepsilon(\mathrm{M}, \mathrm{N})\)，并利用 a) 推出 \(\varepsilon(\mathrm{M}^{\circ}, \mathrm{N}^{\circ}) = \varepsilon(\mathrm{M}, \mathrm{N})\)。
\item
  置 \(g(\mathbf{M}, \mathbf{N}) = \|p_{\mathbf{M}} - p_{\mathbf{N}}\|\)（cf.~V, p.~64, 习题 17）。证明

\[
g (\mathbf {M}, \mathbf {N}) = \sup (\delta (\mathbf {M}, \mathbf {N}), \delta (\mathbf {N}, \mathbf {M}))
\]

（注意 \(p_{\mathbf{M}} - p_{\mathbf{N}} = (1_{\mathrm{E}} - p_{\mathbf{N}})p_{\mathbf{M}} - p_{\mathbf{N}}(1_{\mathrm{E}} - p_{\mathbf{M}})\)）；由此推出 \(\varepsilon (\mathbf{M},\mathbf{N})\leqslant g(\mathbf{M},\mathbf{N})\)。若 \(\mathbf{M}\cap \mathbf{N}^{\circ} = \mathbf{N}\cap \mathbf{M}^{\circ} = \{0\}\)，则我们有关系式 \(\varepsilon (\mathbf{M},\mathbf{N}) = \delta (\mathbf{M},\mathbf{N}) = \delta (\mathbf{N},\mathbf{M}) = g(\mathbf{M},\mathbf{N})\)。若 \(g(\mathbf{M},\mathbf{N}) < 1\)，则 \(\mathbf{M}\cap \mathbf{N}^{\circ} = \mathbf{N}\cap \mathbf{M}^{\circ} = \{0\}\)。

\end{enumerate}

\begin{enumerate}
\def\labelenumi{\alph{enumi})}
\setcounter{enumi}{5}
\tightlist
\item
  设 \(Q_{M}\)、\(Q_{N}\) 是 E 中两个连续投影，它们的像分别为 M 与 N；给出关系式 \(g(\mathrm{M}, \mathrm{N}) \leqslant \|Q_{\mathrm{M}} - Q_{\mathrm{N}}\|\) 的另一个证明（V, p.~64, exerc. 17）。（注意，对所有 \(x \in E\)，我们有 \(\|(1_{E} - Q_{M}).x\|^{2} + \|Q_{M}^{*}.x\|^{2} = \|x\|^{2} + \|(Q_{\mathrm{M}} - Q_{\mathrm{M}}^{*}).x\|^{2}\)，把这一关系式应用于 \(x = (p_{\mathrm{M}} - p_{\mathrm{N}}).y\)，并注意关系式 \((1_{\mathrm{E}} - Q_{\mathrm{M}})(p_{\mathrm{M}} - p_{\mathrm{N}}) = (Q_{\mathrm{M}} - Q_{\mathrm{N}}) p_{\mathrm{N}}\) 与 \((p_{\mathrm{M}} - p_{\mathrm{N}}) Q_{\mathrm{M}} = (1_{\mathrm{E}} - p_{\mathrm{N}})(Q_{\mathrm{M}} - Q_{\mathrm{N}}).\)）
\end{enumerate}

¶ 7) a) 沿用习题 6 的记号，证明：为使 \(\mathbf{M} + \mathbf{N}^{\circ}\) 是闭的，必要且充分的是 \(\mathbf{M} + \mathbf{N}^{\circ} = (\mathbf{M}^{\circ} \cap \mathbf{N})^{\circ}\)。

\begin{enumerate}
\def\labelenumi{\alph{enumi})}
\setcounter{enumi}{1}
\tightlist
\item
  证明下列性质等价：
\end{enumerate}

\(\alpha)\)

\[
\varepsilon (\mathbf {M}, \mathbf {N}) <   1;
\]

β) \(M + N^{\circ}\) 在 E 中是闭的；

\(\gamma)\) 若 \(\tilde{\mathbf{M}}\) 是 \(\mathbf{M} \cap \mathbf{N}^{\circ}\) 在 \(\mathbf{M}\) 中的正交补，而 \((\mathbf{N}^{\circ})^{\sim}\) 是 \(\mathbf{M} \cap \mathbf{N}^{\circ}\) 在 \(\mathbf{N}^{\circ}\) 中的正交补，则 \(\mathbf{E}\) 是 \(\mathbf{M}^{\circ} \cap \mathbf{N}\)、\(\mathbf{M} \cap \mathbf{N}^{\circ}\)、\(\tilde{\mathbf{M}}\) 与 \((\mathbf{N}^{\circ})^{\sim}\) 的直和。此外，若 \(R\) 与 \(S\) 分别是这一分解所对应的、从 \(\mathbf{E}\) 到 \(\mathbf{M}\) 与 \((\mathbf{N}^{\circ})^{\sim}\) 上的投影，则有 \(\| R \| = \| S \| = (1 - \varepsilon^2 (\mathbf{M}, \mathbf{N}))^{1/2}\)。（为看出 \(\alpha\) ) 蕴涵 \(\beta\) )，首先注意：若 \(x \in \tilde{\mathbf{M}}\) 且 \(y \in (\mathrm{N}^\circ)^{\sim}\)，则 \(|\langle x | y \rangle| \leqslant \varepsilon (\mathbf{M}, \mathbf{N}) \| x \| . \| y \|\)；然后，设 \(u = x + y + t\) 是一个元素 \(u \in \mathbf{M} + \mathrm{N}^\circ\) 的分解，其中 \(x \in \mathbf{M}\)、\(y \in (\mathrm{N}^\circ)^{\sim}\) 且 \(t \in \mathbf{M} \cap \mathrm{N}^\circ\)，由此推出 \(\| x \| \leqslant (1 - \varepsilon^2 (\mathbf{M}, \mathbf{N}))^{-1/2} \| u \|, \| y \| \leqslant (1 - \varepsilon^2 (\mathbf{M}, \mathbf{N}))^{-1/2} \| u \|\) 以及 \(\| t \| \leqslant \| u \|\)。为证明 \(\gamma\) ) 蕴涵 \(\alpha\) )，考察分解 \(v = v_1 + v_2\)（对一个 \(v \in \tilde{\mathbf{M}}\)），其中 \(v_1\) 是 \(v\) 到 \(\tilde{\mathbf{N}}\)（即 \(\mathbf{N} \cap \mathbf{M}^\circ\) 在 \(\mathbf{N}\) 中的正交补）上的正交投影；我们有 \(R.v_1 = v\)，且假若有 \(\varepsilon (\mathbf{M}, \mathbf{N}) = 1\)，就会存在一个序列 \((v_n) \in \mathbf{M}\)，使得 \(\| v_n \| = 1\) 且使得 \(\| (v_n)_1\|\) 趋于 0。其次证明限制 \(R_1\)（即 \(R\) 到 \(\tilde{\mathbf{N}}\) 上的限制）是从 \(\tilde{\mathbf{N}}\) 到 \(\tilde{\mathbf{M}}\) 上的一个双射；为计算 \(\| R \|\)，证明 \(\| R_1^{-1} \| \leqslant (1 - \varepsilon^2 (\mathbf{M}, \mathbf{N}))^{1/2}\)。

\begin{enumerate}
\def\labelenumi{\alph{enumi})}
\setcounter{enumi}{2}
\tightlist
\item
  由 b) 推出：若 \(M + N^{\circ}\) 是闭的，则 \(M^{\circ} + N\) 也是闭的。
\end{enumerate}

\begin{enumerate}
\def\labelenumi{\arabic{enumi})}
\setcounter{enumi}{7}
\tightlist
\item
  \begin{enumerate}
  \def\labelenumii{\alph{enumii})}
  \tightlist
  \item
    设 E 是一个希尔伯特空间，T 是一个从 E 到其自身的连续线性映射，满足 \(\|T\|\leqslant1\)。证明关系式 \(T.x=x,\langle T.x|x\rangle=\|x\|\) 与 \(T^{*}.x=x\) 等价，并且 \(1_{E}-T\) 的核与 \(1_{E}-T\) 的像的闭包互为正交补。
  \end{enumerate}
\end{enumerate}

\begin{enumerate}
\def\labelenumi{\alph{enumi})}
\setcounter{enumi}{1}
\tightlist
\item
  设 \(T\) 是一个从 \(E\) 到其自身的连续线性映射，满足不等式

\[
\| x - T. x \| ^ {2} \leqslant \| x \| ^ {2} - \| T. x \| ^ {2}\tag{1}
\]

对所有 \(x \in E\) 成立。于是 \(\|T.x\| < \|x\|\) 对所有使得 \(T.x \neq x\) 的 x 成立；对每个 \(x \in E\)，序列 \((T^{n}.x)\) 收敛于一点 P.x，且 P 是到 \(1_{E} - T\) 的核上的正交投影算子。
\item
  设 \(P_{1}, \ldots, P_{r}\) 是 E 上的正交投影算子。证明乘积 \(T = P_{1}P_{2} \ldots P_{r}\) 满足关系式 (1)（对 r 用归纳法论证）；于是正交投影算子 P 是到诸投影算子 \(P_{j}\) 的像的交上的正交投影算子（注意，若对一个指标 j 有 \(\|P_{j}.x\| < \|x\|\)，则 \(\|T.x\| < \|x\|\)）。
\end{enumerate}

\begin{enumerate}
\def\labelenumi{\arabic{enumi})}
\setcounter{enumi}{8}
\tightlist
\item
  设 E 是一个希尔伯特空间，\((P_{j})_{j\in\mathbb{N}}\) 是 E 上正交投影算子的一个序列，使得对所有 \(j\in N\)，存在一个 \(n_{j}\in N\)，使得对所有 \(k\in N\)，正交投影算子 \(P_{k}, P_{k+1}, ..., P_{k+n_{j}}\) 中至少有一个等于 \(P_{j}\)。置 \(R_{s}=P_{s}P_{s-1}\ldots P_{0}\)（对所有 \(s\in N\)）。
\end{enumerate}

\begin{enumerate}
\def\labelenumi{\alph{enumi})}
\tightlist
\item
  对每个 \(x\in E\)，令 \(x_{s}=R_{s}\cdot x\)。证明 \(\sum_{s}\|x_{s-1}-x_{s}\|^{2}\leqslant\|x\|^{2}\)，并由此推出，对
\end{enumerate}

每个整数 \(r \geqslant 1\)，\(x_{s + r} - x_s\) 当 \(s\) 趋于 \(+\infty\) 时趋于 0。

\begin{enumerate}
\def\labelenumi{\alph{enumi})}
\setcounter{enumi}{1}
\item
  设 \((x_{s_k})\) 是从 \((x_s)\) 中抽取的一个序列，它弱收敛于一个极限 \(y\)；于是每个序列 \((x_{s_k + r})\) 也弱收敛于 \(y\)。由此推出，\(y\) 属于每个子空间 \(\mathbf{M}_j = P_j(\mathbf{E})\)。（对每个 \(j\)，存在 \(r_k\) 使得 \(0 \leqslant r_k \leqslant n_j\) 且 \(s_k + r_k = j\)；证明序列 \((x_{s_k + r_k})\) 弱收敛于 \(y\)。）
\item
  证明序列 \((x_{s})\) 弱收敛于 \(x\) 到交 \(\mathbf{M}\)（诸 \(\mathbf{M}_{j}\) 的交）上的正交投影。（化归到 \(\mathbf{M} = \{0\}\) 的情形，并利用 \(b\) ) 以及 E 中每个闭球的弱紧性。）
\end{enumerate}

¶ 10) 设 E 是一个希尔伯特空间，u 是 E 的一个正算子。

\begin{enumerate}
\def\labelenumi{\alph{enumi})}
\tightlist
\item
  证明：对所有 \(x \in E\)，我们有

\[
\left\| u (x) \right\| ^ {2} \leqslant \| u \|. \langle u (x) | x \rangle
\]

(注意 \(\langle u(x)|u(x)\rangle^2\leqslant \langle u(x)|x\rangle \langle u^2 (x)|u(x)\rangle\)，见 V, p.~3, prop. 2)。
\item
  设 M 是 E 的一个闭向量子空间，\(\mathbf{M}^{\circ}\) 是它的正交补。设 \(x\in \mathbf{M}\)，并令 \(f(x)\) 为 \(\langle u(x + y)|x + y\rangle\)（当 \(y\) 取遍 \(\mathbf{M}^{\circ}\) 时）的下确界。对每个 \(\varepsilon >0\)，令 \(\operatorname {E}(x,\varepsilon)\) 为由所有 \(y\in \mathbf{M}^{\circ}\)（满足 \(\langle u(x + y)|x + y\rangle \leqslant f(x) + \varepsilon\)）组成的集。证明 \(\operatorname {E}(x,\varepsilon)\) 是凸的，并且对所有 \(z\in \mathbf{M}^{\circ}\)，我们有

\[
\langle u (x + y) | z \rangle^ {2} \leqslant \varepsilon \langle u (z) | z \rangle \quad \text { for } \quad y \in \mathrm{E} (x, \varepsilon)\tag{*}
\]

(考察实变量 t 的函数 \(g:t \mapsto \langle u(x + y + tz)/x + y + tz \rangle\)，它在一点 \(t_{0}\) 处达到其极小值，并注意 \(g(t_{0}) \geqslant f(x)\) 与 \(g(0) \leqslant f(x) + \varepsilon\) )。
\item
  对每个整数 \(n \geqslant 1\)，取 \(y_{n} \in \mathrm{E}(x, 1/n)\)；证明序列 \((u(x + y_{n}))\) 趋于一个属于 M 的极限 \(x_{1}\)，且序列 \((\langle u(y_{n})|y_{n}\rangle)\) 有界（利用不等式 (*)，为数 \(\langle u(y_{n} - y_{m})|y_{n} - y_{m}\rangle\)（对 \(m \geqslant n\)）以及数 \(|\langle u(x + y_{n})|z\rangle|\)（对 \(z \in M^{\circ}\)）各找出一个界）。

\item
  设 \((y_n')\) 是 \(\mathbf{M}^\circ\) 中点的一个序列，使得 \(\langle u(y_n' - y_m')|y_n' - y_m'\rangle\) 当 \(m\) 与 \(n\) 充分大时任意小，且使得序列 \((u(x + y_n'))\) 有一个极限 \(x_1' \in \mathbf{M}\)；

证明 \(x_1' = x_1\)。（令 \(Q(z) = \langle u(z)|z \rangle\)；首先证明数 \(Q((y_p - y_p') - (y_q - y_q'))\) 一旦 \(p\) 与 \(q\) 充分大就任意小，并由此推出序列 \((Q(y_n - y_n'))\) 有界；利用 \(\langle x_1' - x_1|y_p - y_p' \rangle = 0\)（对所有 \(p\)）这一事实，证明序列 \((Q(y_n - y_n'))\) 趋于 0，并利用 \(a\) )。

\item
  由 \(d\) ) 推出：点 \(x_{1}\) 不依赖于 \(y_{n} \in \mathrm{E}(x, 1/n)\) 的选取，且若置 \(u_{1}(x) = x_{1}\)，则 \(u_{1}\) 是一个从 \(\mathbf{M}\) 到其自身的线性映射。证明 \(0 \leqslant \langle u_{1}(x)|x\rangle \leqslant \langle u(x)|x\rangle\) 对所有 \(x \in \mathbf{M}\) 成立，因而 \(u_{1}\) 连续，并且是 \(\mathbf{M}\) 的一个 \(\geqslant 0\) 的自同态。（注意 \(\langle u(x + y_n)|y_n\rangle\) 趋于 0 且 \(\langle u(x + y_n)|x + y_n\rangle\) 趋于 \(f(x)\)。）
\item
  设 \(p_{M}\) 是以 M 为像的正交投影算子，并令 \(u_{0} = u_{1} \circ p_{M}\)。我们有 \(0 \leqslant u_{0} \leqslant u\)，M 在 \(u_{0}\) 下稳定，且 \(u_{0}\) 到 \(M^{\circ}\) 上的限制为零。证明 \(u_{0}\) 是由所有满足 \(v \geqslant 0\)、\(v \leqslant u\)、M 在 v 下稳定且 v 到 \(M^{\circ}\) 上的限制为零的自同态 v 组成的族中的最大元素。
\end{enumerate}

\begin{enumerate}
\def\labelenumi{\arabic{enumi})}
\setcounter{enumi}{10}
\tightlist
\item
  设 \(\mathrm{E}\) 与 \(\mathrm{F}\) 是两个希尔伯特空间。证明：对每个元素 \(u\)（\(\mathrm{E} \hat{\otimes}_2 \mathrm{F}\) 中的），存在一个标准正交序列 \((e_n)\)（\(\mathrm{E}\) 中的）、一个标准正交序列 \((f_n)\)（\(\mathrm{F}\) 中的）以及一个序列 \((\lambda_n)\)（其诸数 \(\geqslant 0\)），使得 \(\sum_{n} \lambda_n^2 < +\infty\) 且 \(u = \sum_{n} \lambda_n e_n \otimes f_n\)；于是 \(\| u \|_2^2 = \sum_{n} \lambda_n^2\)
\end{enumerate}

(cf.~V, p.~55, th. 2 及 p.~53, th. 1)。

¶ 12) 设 E 是一个实希尔伯特空间，V 是 E 中一个以 0 为顶点的闭凸锥，令 V° 为 V 的极锥（在 E 中，把 E 典范地与其对偶等同）。

\begin{enumerate}
\def\labelenumi{\alph{enumi})}
\item
  证明每个点 \(x \in \mathbf{E}\) 都可以唯一地写成形式 \(x = x_{+} - x_{-}\)，其中 \(x_{+} \in \mathbf{V}\) 且 \(x_{-} \in \mathbf{V}^{\circ}\)，以及 \(\langle x_{+}|x_{-}\rangle = 0\)。
\item
  对 V 的每个刻面 F（II, p.~87, exerc. 3），F 或是点 0，或是一个以 0 为顶点的凸锥；所有与 F 正交的 \(y \in V^{\circ}\) 组成的集是一个闭刻面 \(F'\)（\(V^{\circ}\) 的）（但它不是 II, p.~87, exerc. 6 意义下 F 的「对偶刻面」，后者是空集）。
\item
  取 E 为一个实希尔伯特空间的全体 Hilbert-Schmidt 自同态组成的集，取 V 为 E 中正元素组成的集。证明我们有 \(V^{\circ} = V\)，并在此情形解释 a) 的结果（为看出 \(V \subset V^{\circ}\)，利用 V, p.~56 的 cor. 1）。
\item
  在 c) 的假设下，设 \(v \in V\)；所有 \(x \in H\)（满足 \(\langle v(x)|x \rangle = 0\)，或等价地 \(v(x) = 0\)，V, p.~77, exerc. 10）组成的集 L 是 H 的一个闭向量子空间，且 v 在 V 中的刻面 F 是由所有这样的 \(u \in V\) 组成的闭集：\(u(x) = 0\) 对所有 \(x \in L\) 成立；它可以与希尔伯特空间 \(L^{\circ}\) 的全体正 Hilbert-Schmidt 自同态组成的锥等同。由此推出：从 E 到凸集 F 上的投影（V, p.~11）与从 E 到由 F 生成的 E 的闭向量子空间上的正交投影算子恒同。
\end{enumerate}

¶ 13) 设 E 是一个希尔伯特空间，G 是 E 的希尔伯特空间结构的全体自同构组成的群的一个子群。令 \(E^{G}\) 为 E 中由在 G 下不变的所有向量组成的闭向量子空间，p 为从 E 到 \(E^{G}\) 上的正交投影算子。

\begin{enumerate}
\def\labelenumi{\alph{enumi})}
\item
  证明 \(\mathbf{E}^{\mathrm{G}}\) 在 \(\mathbf{E}\) 中的正交补是由向量 \(s.x - x\)（其中 \(s\) 取遍 \(\mathbf{G}\) 且 \(x \in \mathbf{E}\)）生成的闭向量子空间。
\item
  设 \(\mathbf{H}\) 是 \(\mathbf{E}\) 的一个在 \(\mathbf{G}\) 下稳定的非空闭凸子集。证明 0 到 \(\mathbf{H}\) 上的投影属于 \(\mathbf{E}^{\mathrm{G}}\)。
\item
  假设 H 是 E 中一点 x 的轨道的闭凸包，令 a 为 0 到 H 上的投影。证明 \(a = p(x)\)，且 \(H \cap E^{G}\) 缩为点 a（「Birkhoff-Alaoglu 定理」）。（注意 x - a 包含在 \(E^{G}\) 的正交补中。）
\item
  假设 G 由 E 的一个自同构 \(u\) 生成。证明 \(p(x) = \lim_{n\to \infty}\frac{1}{n + 1}\sum_{j = 0}^{n}u^{j}(x)\)
\end{enumerate}

对所有 \(x \in \mathrm{E}\) 成立（若 \(y_{n} = \frac{1}{n + 1} \sum_{j=0}^{n} u^{j}(x)\)，注意序列 \((y_{n})\) 有一个弱极限点 \(a\)，且 \(u(a) = a\)，然后利用 \(c\) )。

\begin{enumerate}
\def\labelenumi{\alph{enumi})}
\setcounter{enumi}{4}
\tightlist
\item
  假设 \(G\) 是一个同态 \(t \mapsto u_t\)（从 \(\mathbf{R}\) 到 \(E\) 的自同构群上的）的像，使得对所有 \(x \in E\)，\(t \mapsto u_t \cdot x\) 是一个从 \(\mathbf{R}\) 到 \(E\) 中的连续映射。证明

\(p(x) = \lim_{T \to \infty} \frac{1}{T} \int_{0}^{T} u_t \cdot x \, dt\)

对所有 \(x \in E\) 成立。

\end{enumerate}

\begin{enumerate}
\def\labelenumi{\alph{enumi})}
\setcounter{enumi}{5}
\tightlist
\item
  假设存在一个元素 \(x \neq 0\)（属于 \(\mathbf{E}\)）和一个数 \(\alpha\)，使得 \(0 < \alpha < 1\) 且 \(\| s.x - x\| \leqslant \alpha \| x\|\) 对所有 \(s \in \mathbf{G}\) 成立。证明 \(\mathrm{E}^{\mathrm{G}} \neq \{0\}\)（利用 \(c\)）。
\end{enumerate}

\begin{enumerate}
\def\labelenumi{\arabic{enumi})}
\setcounter{enumi}{13}
\tightlist
\item
  设 E 是一个复希尔伯特空间，T 是 E 的一个 Hilbert-Schmidt 自同态。a) 设 R 与 L 是正 Hilbert-Schmidt 自同态，满足 \(R^{2} = T^{*}T\) 与 \(L^{2} = TT^{*}\)（V, p.~57, cor. 3）；令 \(R = \text{abs}(T)\)，并称之为 T 的「绝对值」（cf.~V, p.~75, exerc. 4）；我们有 \(L = \text{abs}(T^{*})\)。证明 \(\text{Ker}(T) = \text{Ker}(R)\) 且 \(\overline{L(\text{E})} = \overline{T(\text{E})}\)。存在唯一的等距 V，它从 \(R(\text{E})\) 到 \(T(\text{E})\) 上，使得 T = VR；若把 V 连续地延拓到 \(\overline{R(\text{E})}\)，再延拓为一个算子 \(U \in \mathcal{L}(\text{E})\)（通过在 \(R(\text{E})\) 的正交补上取 U.x = 0），则我们仍有 T = UR（T 的极分解）。于是 \(R = U^{*}T = U^{*}UR = RU^{*}U\) 且 \(L = URU^{*}\)、\(T = LU^{*}\)。若 T 属于 \(\mathcal{L}^{1}(\text{E})\)，则 \(R = \text{abs}(T)\) 亦然，且 T 是两个 Hilbert-Schmidt 自同态的乘积。
\end{enumerate}

\begin{enumerate}
\def\labelenumi{\alph{enumi})}
\setcounter{enumi}{1}
\item
  若 \(T\) 属于 \(\mathcal{L}^1 (\mathrm{E})\)，证明 \(\operatorname {Tr}(\operatorname {abs}(T)) = \sup (\sum_i |\langle a_i | T.b_i\rangle |)\)，其中在右端，\((a_{i})\) 与 \((b_{i})\) 取遍 E 的标准正交基组成的集（利用 \(T\) 的极分解）。证明：若置 \(\| T\| _1 = \operatorname {Tr}(\operatorname {abs}(T))\)，则 \(\| T\| _1\) 是空间 \(\mathcal{L}^1 (\mathrm{E})\) 上的一个范数，使得 \(\| T\| _2\leqslant \| T\| _1\)。
\item
  反之，若 \(T \in \mathcal{L}(\mathrm{E})\) 使得对每一对 \(((a_i), (b_i))\)（\(\mathbf{E}\) 的标准正交基的），和式 \(\sum_{i} |\langle a_i | T \cdot b_i \rangle|\) 都有限，则 \(T \in \mathcal{L}^1(\mathrm{E})\)（首先注意 \(T\) 是一个 Hilbert-Schmidt 自同态，然后利用 \(T\) 的极分解）。
\item
  设 \((\bar{T}_{\nu})\) 是 Hilbert-Schmidt 自同态（相应地 \(\mathcal{L}^1 (\mathrm{E})\) 的元素）的一个序列，使得对 E 的每对点 \(x,y\)，序列 \((\langle x|T_{\nu}.y\rangle)\) 收敛于 \(\langle x|T.y\rangle\)，其中 \(T\) 是一个从 E 到其自身的线性映射；此外，假设范数序列 \(\| T_{\nu}\| _2\)（相应地 \(\| T_{\nu}\| _1\)）有界。证明 \(T\) 是一个 Hilbert-Schmidt 自同态（相应地 \(\mathcal{L}^1 (\mathrm{E})\) 的一个元素）（利用 \(b\) ) 与 \(c\) )）。
\item
  由 \(d\) ) 推出：空间 \(\mathcal{L}^1 (\mathrm{E})\) 对范数 \(\| T\| _1\) 是一个 Banach 空间。
\item
  为使一个自同态 \(T \in \mathcal{L}(\mathrm{E})\) 属于 \(\mathcal{L}^1(\mathrm{E})\)，必要且充分的是：对至少一个标准正交基 \((e_i)\)（\(\mathrm{E}\) 的），和式 \(\sum_{i} \| T.e_i\|\) 有限（沿用 \(a\) ) 的记号，注意 \(|\langle e_i|R.e_i\rangle| \leqslant \| T.e_i\|\)）。

\item
  在空间 \(\ell_{\mathbf{C}}^2\) 中，令 \((e_n)\) 为典范标准正交基，并令 \(a = \sum_{n=0}^{\infty} \frac{1}{n+1} e_n\)；

若 \(\mathrm{F}\) 是一维子空间 \(\mathbf{C}.a\)，则正交投影算子 \(p_{\mathrm{F}}\) 有有限迹，但级数 \(\sum_{n}\| p_{\mathrm{F}}\cdot e_n\|\) 不收敛。
\end{enumerate}

\begin{enumerate}
\def\labelenumi{\arabic{enumi})}
\setcounter{enumi}{14}
\tightlist
\item
  设 \(\mathbf{E}\) 是一个复希尔伯特空间；令 \(\mathcal{B} = \mathcal{L}(\mathbf{E})\) 表示 \(\mathbf{E}\) 的连续自同态组成的代数，赋以通常的范数 \(\| T\| = \sup_{\| x\| \leqslant 1}\| T.x\|\)。对 E 的每对点 \(x, y\)，令 \(\omega_{x,y}\) 表示连续线性型 \(T \mapsto \langle x|T.y \rangle\)（\(\mathcal{B}\) 上的），并令 \(\mathcal{B}_{\circ}\) 为强对偶 \(\mathcal{B}'\)（Banach 空间 \(\mathcal{B}\) 的）中由诸 \(\omega_{x,y}\) 生成的闭子空间。a) 证明：把每个 \(T \in \mathcal{B}\) 对应到线性型 \(\omega \mapsto \langle \omega, T \rangle\)（\(\mathcal{B}_{\circ}\) 上的）的线性映射，是从 \(\mathcal{B}\) 到 \(\mathcal{B}_{\circ}\) 的强对偶上的一个等距；换句话说，\(\mathcal{B}_{\circ}\) 是 \(\mathcal{B}\) 的一个预对偶（IV, p.~56, exerc. 23）。
\end{enumerate}

拓扑 \(\sigma (\mathcal{B},\mathcal{B}_{\circ})\)（\(\mathcal{B}\) 上的）称为超弱拓扑。

\begin{enumerate}
\def\labelenumi{\alph{enumi})}
\setcounter{enumi}{1}
\item
  对每个元素 \(T\)（\(\mathcal{L}^1(\mathrm{E})\) 中的），我们定义一个线性型 \(\phi_T\)（\(\mathcal{B}\) 上的），所用公式为 \(\phi_T(S) = \operatorname{Tr}(ST)\)（对每个算子 \(S \in \mathcal{B}\)）。证明 \(\phi_T\) 连续，且映射 \(T \mapsto \phi_T\) 是从 Banach 空间 \(\mathcal{L}^1(\mathrm{E})\)（exerc. 14, e)）到 Banach 空间 \(\mathcal{B}_{\circ}\) 上的一个等距（先考虑 \(T\) 为有限秩的情形）。
\item
  令 \(\mathcal{B}_{\circ \circ}\) 为 \(\mathcal{B}_{\circ}\) 中由诸 \(\omega_{x,y}\) 生成的向量子空间（于是 \(\mathcal{B}_{\circ} = \overline{\mathcal{B}}_{\circ \circ} \)）。证明 \(\mathcal{B}_{\circ \circ}\) 是桶型的（注意 \(\mathcal{B}\) 的一个子集若对 \(\sigma (\mathcal{B},\mathcal{B}_{\circ \circ})\) 有界，则对范数拓扑有界）。
\item
  令 \(F_{n}\) 为 \(\mathcal{B}_{\circ \circ}\) 的一个子空间，它是 E 的秩 \(\leqslant n\) 的自同态组成的集在 b) 中定义的等距下的像。证明 \(F_{n}\) 在 \(\mathcal{B}_{\circ \circ}\) 中无处稠密，并由此推出 \(\mathcal{B}_{\circ \circ}\) 不是 Baire 空间。
\end{enumerate}

\backmatter
\backmatterentry{历史注记}

(第 I 至 V 章)

(N.B. --- 罗马字母指本注记末尾的参考文献。)

拓扑向量空间的一般理论创立于 1920 至 1930 年前后的时期。但其基础工作早已通过对泛函分析众多问题的研究而长期准备着；若不至少简略地指出这些问题的研究如何缓慢地（特别自 20 世纪初以来）使数学家们认识到所考虑的各种问题之间的联系，以及以远为一般的方式表述这些问题、并对它们应用统一的求解方法的可能性，我们就无法追溯这一学科的历史。

可以说，代数与分析之间的类比，以及把函数方程（即未知量为函数的方程）视为代数方程的「极限情形」的想法，都起源于无穷小演算，后者在某种意义上正是为了「从有限到无穷」进行推广而发明的。但无穷小演算的直接代数祖先是有限差分演算（cf.~FVR, 第 I、II、III 章的历史注记，p.~54-58），而不是一般线性方程组的求解；直到 18 世纪中叶以后，后者与微分演算问题之间的最初类比才在弦振动方程的研究中出现。我们在此不打算深入讨论这一问题的历史细节；但有两个基本想法的不断重现令人瞩目，它们显然都应归功于 D. Bernoulli。第一个想法在于把弦的振荡视为 n 个质点组成的系统的振荡当 n 无限增加时的「极限情形」；我们知道，后来这个 n 为有限的问题成了寻求线性变换特征值的第一个例子（cf.~A, 第 VI-VII 章的历史注记）；这些数在极限情形对应于弦的「本征振荡」的频率，后者早在很久以前就被实验观察到，而其理论上的存在性在该世纪初（特别由 Taylor）已被确立。这种形式上的类比，虽然在后来几乎从未被提及（(I, b), p.~390），但在整个 19 世纪中似乎从未被忽视；但正如我们将看到的，只是到了 1890-1900 年前后，它才获得了全部的重要性。

D. Bernoulli 的另一个想法（或许受实验事实启发）是「叠加原理」，按照这一原理，弦的最一般的振荡应当可以通过「本征振荡」的叠加被「分解」；用数学的语言来说，这意味着弦振动方程的通解应当有形如 \(\sum_{n} c_{n} \phi_{n}(x, t)\) 的级数展开，其中 \(\phi_{n}(x, t)\)

表示本征振荡。我们知道，这一原理是关于「任意」函数能否展开为三角级数的一场长期论战的出发点，这场论战直到 19 世纪头三分之一才由 Fourier 与 Dirichlet 的工作解决。但甚至在这一结果获得之前，就已有其他按「正交」函数 * 展开级数的例子：球面函数、Legendre 多项式，以及形如 ( \(e^{i\lambda_n x}\) ) 的各种函数系，其中 \(\lambda_n\) 不再是同一个数的倍数；这些函数系在 18 世纪就已引入振荡问题之中，同样也由 Fourier 与 Poisson 在他们关于热理论的研究中引入。1830 年前后，Sturm (I) 与 Liouville (II) 把在这些各种特殊情形中观察到的全部现象系统化，建立起一元函数振荡的一般理论；他们考虑微分方程

\[
\frac {d}{d x} \left(p (x) \frac {d y}{d x}\right) + \lambda \rho (x) y = 0 \quad (p (x) > 0, \rho (x) > 0)\tag{1}
\]

以及边界条件

\[
\begin{array}{l l} k _ {1} y ^ {\prime} (a) - h _ {1} y (a) = 0 \\ k _ {2} y ^ {\prime} (b) + h _ {2} y (b) = 0 \end{array} (h _ {1} k _ {1} \neq 0, h _ {2} k _ {2} \neq 0, a <   b)\tag{2}
\]

并证明了下列基本结果：

\begin{enumerate}
\def\labelenumi{\arabic{enumi})}
\item
  仅当 \(\lambda\) 取一个序列 \((\lambda_{n})\)（由数 > 0 组成且趋于 \(+\infty\)）中的值之一时，此问题才有非零解；
\item
  对每个 \(\lambda_{n}\)，解都是同一个函数 \(v_{n}\) 的倍数，后者可以借助条件 \(\int_{a}^{b} \rho v_{n}^{2} dx = 1\) 假定为「规范化的」，且对 \(m \neq n\) 我们有 \(\int_{a}^{b} \rho v_{m} v_{n} dx = 0\)；
\item
  每个函数 \(f\)（在 \([a, b]\) 上二次可微且满足边界条件 (2) 的），都可以展开为一致收敛的级数 \(f(x) = \sum_{n} c_{n} v_{n}(x)\)，
\end{enumerate}

其中 \(c_{n} = \int_{a}^{b}\rho fv_{n}dx\)

\begin{enumerate}
\def\labelenumi{\arabic{enumi})}
\setcounter{enumi}{3}
\tightlist
\item
  等式 \(\int_{a}^{b}\rho f^2 dx = \sum_n c_n^2\) 成立（这一等式早在 1799 年就已被 Parseval 以纯形式的方式对三角函数系证明；由它立即可推出「Bessel 不等式」；后一不等式由 Bessel 于 1828 年（仍就三角级数）公布）。
\end{enumerate}

半个世纪之后，这些性质由 Gram (III) 的工作补充完备，他在 Tchebichef 的研究之后，阐明了正交函数级数展开与「最佳平方逼近」问题（Gauss 误差论中「最小二乘法」的直接产物）之间的关系；后者如下：给定一个函数的有穷序列 \((\psi_{i})_{1\leqslant i\leqslant n}\) 和一个函数 f，求线性组合 \(\sum_{i}a_{i}\psi_{i}\)，使得积分 \(\int_{a}^{b}\rho (f - \sum_{i}a_{i}\psi_{i})^{2}dx\) 达到其极小值。原则上，这只是一个平凡的线性代数问题，但 Gram 以一种独创的方式解决了它，即对 \(\psi_{i}\) 应用如 chap.~V, p.~23 中所述的（且通常以 Erhard Schmidt 之名为人所知的）「标准正交化」方法。其次，在无穷标准正交系的情形，产生的问题是：一个函数 f 借助于序列前 n 个函数的线性组合所得的「最佳平方逼近」\(\mu_{n}\)，何时在 n 无限增加时趋于 0 *；Gram 由此被引向完全标准正交系这一概念的定义，并认识到这一性质等价于不存在与所有 \(\phi_{n}\) 正交的非零函数。他甚至试图阐明「平均平方收敛」的概念，但在测度论的基本思想被引入之前，他在这一方向上很难得到任何一般的结果。
* 但这一术语在 Hilbert 的工作之前并未出现。

在 19 世纪的下半叶，分析学家们的主要力量主要用于把 Sturm-Liouville 理论推广到多元函数。这一理论是由数学物理中出现的椭圆型偏微分方程以及自然伴随它们的边值问题所促动的。主要的兴趣首先集中在「振动膜」方程上

\[
\mathrm{L} _ {\lambda} (u) \equiv \Delta u + \lambda u = 0\tag{3}
\]

其中所求的是在一个充分正则的区域 G 的边界上为零的解；对一元函数行之有效的方法对这一问题不再适用，出现的巨大分析困难被逐步克服。我们回顾通往解的主要步骤：G 的「Green 函数」的引入，其存在性由 Schwarz 证明；最小特征值的存在性的证明，同样出自 Schwarz；最后，1894 年，H. Poincaré 在一篇著名的论文 (V a) 中成功地证明了所有特征值的存在性及其基本性质。他对给定的「右端」f 考虑方程 \(\mathbf{L}_{\lambda}(u) = f\) 的解；该解在边界上为零；然后，通过对 Schwarz 方法的一个巧妙推广，他证明了 \(u_{\lambda}\) 是复变量 \(\lambda\) 的亚纯函数，只有实的单极点 \(\lambda_{n}\)，而这些正是所求的特征值。

* 必须指出，在这一研究中，Gram 并不把自己局限于只考虑连续函数，而是强调了条件 \(\int_{a}^{b} \rho f^2 dx < +\infty\) 的重要性。

这些研究与线性积分方程论的开端直接相关，后者必定对现代思想的到来做出了最大的贡献。我们在此只限于给出这一理论发展的一个简略概要（更完整的细节，请参看本专著后面专讲谱论的各章之后的历史注记）。这类函数方程最初在 19 世纪上半叶不显眼地出现（Abel、Liouville），自从 Beer 与 C. Neumann 把一个充分正则的区域 G 上的「Dirichlet 问题」的解化为对一个「第二类积分方程」

\[
u (x) + \int_ {a} ^ {b} \mathrm{K} (x, y) u (y) d y = f (x)\tag{4}
\]

(对未知函数 u)的求解以来，它就已经获得了一定的重要性；C. Neumann 于 1877 年用「逐次逼近」法成功地解出了这个方程。H. Poincaré 一方面受到上述代数类比的启发，另一方面受到他为振动膜方程所得结果的启发，于 1896 年 (V b) 在前一方程的积分前引入了一个变参数 \(\lambda\)，并断言，正如在振动膜方程的情形那样，解仍然是 \(\lambda\) 的亚纯函数；但他未能证明这一结果。七年后，I. Fredholm (VI) 建立了它（对连续的「核」K 与有穷区间 \([a, b]\)）。刚才提到的这位作者，或许比他的前辈们有更清醒的意识，任凭自己由 (4) 与线性方程组

\[
\sum_ {q = 1} ^ {n} \left(\delta_ {p q} + \frac {1}{n} a _ {p q}\right) x _ {q} = b _ {p} (1 \leqslant p \leqslant n)\tag{5}
\]

的类比引导，从而把 (4) 的解表为两个表达式之商，这两个表达式仿照行列式的模式，正如 Cramer 公式中出现的那些。然而这并不是一个新想法：自 19 世纪初以来，「待定系数」法（其做法是：假定一个未知函数有级数展开 \(\sum_{n} c_{n} \phi_{n}\)，其中 \(\phi_{n}\) 是已知函数，通过计算系数 \(c_{n}\) 来求出这个未知函数）就已引出「具有无穷多个未知量的线性方程组」

\[
\sum_ {j = 1} ^ {\infty} a _ {i j} x _ {j} = b _ {i} (i = 1, 2, \dots).\tag{6}
\]

Fourier 遇到过这样一个方程组，他仍像 18 世纪的数学家那样求解它：他删去所有指标 \(i\) 或 \(j\) 大于 \(n\) 的项，用 Cramer 公式显式地解出如此得到的有穷方程组，然后在解中令 \(n\) 趋于 \(+\infty\) 取极限！很久以后，当这套戏法不再能被接受时，问题仍是用行列式论来处理的；自 1886 年起（继 Hill 的工作之后），H. Poincaré 以及后来的 H. von Koch 建立了一种「无穷行列式」理论，它使得某些类型的方程组 (6) 可以按经典模式求解；而且尽管这些结果并不能直接适用于 Fredholm 所处理的问题，von Koch 的理论尤其无疑曾充当了构造 Fredholm「行列式」的模型。

正是在这时，Hilbert 登场并给这一理论以新的推动 (VII)。首先，他完成了 Fredholm 的工作，切实地实现了从 (5) 的解到 (4) 的解的极限过渡；但他随即引入了实二次型理论中相应的极限过渡，这种二次型从带对称核（即满足 \(\mathbf{K}(y,x)=\mathbf{K}(x,y)\)）的积分方程中自动产生。这些方程是数学物理中最常见的方程。他由此成功地得到了直接推广二次型化到主轴的基本公式

\[
\int_ {a} ^ {b} \int_ {a} ^ {b} \mathrm{K} (s, t) x (s) x (t) d s d t = \sum_ {n = 1} ^ {\infty} \frac {1}{\lambda_ {n}} \left(\int_ {a} ^ {b} \phi_ {n} (s) x (s) d s\right) ^ {2},\tag{7}
\]

其中 \(\lambda_{n}\) 是核 K 的特征值（必然是实的），\(\phi_{n}\) 构成相应本征函数的标准正交系，且当 \(\int_{a}^{b} x^{2}(s) ds \leqslant 1\) 时，公式 (7) 的右端是一个收敛级数。他还指出了每个「可表示」为 \(f(x) = \int_{a}^{b} \mathbf{K}(x, y) g(y) dy\) 的函数如何具有一个「展开」 \(\sum_{n=1}^{\infty} \phi_{n}(x) \int_{a}^{b} \phi_{n}(y) f(y) dy\)，并且，依照与二次型经典理论的类比，他给出了一个用变分方法确定 \(\lambda_{n}\) 的程序。这正是二次曲面主轴的那些著名极值性质的推广 ((VII), p.~1-38)。

Hilbert 的这些初步结果几乎立即被 E. Schmidt 以更简单且更一般的形式接手，避免了「Fredholm 行列式」的引入，也避免了从有限到无穷的过渡。其表述已非常接近抽象，积分的线性性与正性这两个基本性质显然是证明中所用到的仅有的事实 (VII \(a\))。但到那时，Hilbert 已发展了远为一般的概念。所有更早的工作都显示出平方可积函数的重要性，而 Parseval 公式在这些函数与序列 ( \(c_n\) )（满足 \(\sum_{n} c_n^2 < \infty\)）之间建立了直接联系。正是这一想法引导着 Hilbert 在他 1906 年的论文中 ((VII), chap.~XI-XIII)，再次采用「待定系数」的旧方法，证明了积分方程 (4) 的求解等价于一个无穷线性方程组的求解

\[
x _ {p} + \sum_ {q = 1} ^ {\infty} k _ {p q} x _ {q} = b _ {p} (p = 1, 2, \dots)\tag{8}
\]

「Fourier 系数」\(x_{p} = \int_{a}^{b} u(t) \omega_{p}(t) dt\)（未知函数 \(u\) 关于一个给定的完全标准正交系 \((\omega_{n})\) 的，其中 \(b_{p} = \int_{a}^{b} f(t) \omega_{p}(t) dt\)，\(k_{pq} = \int_{a}^{b} \int_{a}^{b} K(s, t) \omega_{p}(s) \omega_{q}(t) ds dt\)）。此外，从这一观点看，(8) 的解中唯一令人感兴趣的是使 \(\sum_{n} x_{n}^{2} < +\infty\) 的那些解；Hilbert 也正是系统地把自己局限于这类解；但另一方面，他放宽了加于「无穷矩阵」\(k_{pq}\) 上的条件（它在 (8) 中满足 \(\sum_{p,q} k_{pq}^{2} < +\infty\)）。从那时起，情况已经清楚：由所有实数序列 \(x = (x_{n})\)（满足 \(\sum_{n} x_{n}^{2} < +\infty\)）组成的「Hilbert 空间」，虽然没有被明确引入，却是整个理论底层的基础空间，并且表现为有限维欧氏空间的一种「极限过渡」。此外——这对后来的发展特别重要——Hilbert 被引着在这个空间中引入的不是一个、而是两个不同的收敛概念（相应于后来所谓的弱拓扑与强拓扑 *），同样还有一个「选择原理」，它恰好是单位球的弱紧性。他为解方程组 (8) 而发展的新线性代数完全依赖这些拓扑思想：线性映射、线性型与双线性型（与线性映射相伴的）按照它们的「连续性」性质 ** 被分类和研究。特别地，Hilbert 发现 Fredholm 方法的成功依赖于「完全连续」概念，这一概念通过他对双线性型 *** 表述它并深入加以研究而得到挽救；更详细的细节，我们请读者参看本专著中将要展开这一重要概念的部分，以及 Hilbert 那些令人赞叹的深刻著作，他在其中开创了对称双线性型（有界或非有界）的谱理论。

Hilbert 的语言仍是经典的语言，而且在整个「Grundzüge」中，他从未忽视他所发展的理论的应用，这些应用取自大量例子（几乎占了全书一半篇幅）。下一代人已经采取了抽象得多的观点。在 Fréchet 与 F. Riesz 关于一般拓扑的思想的影响下（见 GT, chap.~I 的历史注记），E. Schmidt (VII b) 与 Fréchet 本人在 1907-1908 年间自觉地把欧氏几何的语言引入「Hilbert 空间」（实的或复的）；正是在这些著作中，我们首次见到范数（连同它现在的记号 \(\| x \|\)）、它所满足的三角不等式，以及 Hilbert 空间是「可分的」且完备这一事实；此外，E. Schmidt 证明了到闭线性流形上的正交投影的存在性，这使他能够把 Hilbert 的线性方程组理论表述成更简单且更一般的形式。同样在 1907 年，Fréchet 与 F. Riesz 注意到平方可积函数空间有一种类似的「几何」，几个月后，当 F. Riesz 与 E. Fischer 证明了这个空间是完备的且同构于一个「Hilbert 空间」时，这一类比得到了完美的解释，同时以引人注目的方式显示了 Lebesgue 新创造的工具的价值。从这时起，希尔伯特空间理论的要点就可以被认为已经完成。在后来的发展中，必须提到 M. H. Stone 与 J. von Neumann 在 1930 年前后对该理论的公理化表述，以及「可分性」限制的取消，后者是 Rellich、Löwig 与 F. Riesz (IX e) 在 1934 年前后工作的成果。
* 变分法自然地导致了同一函数集上不同的收敛概念（依函数一致收敛的要求，或依函数及其若干导数一致收敛的要求）；但 Hilbert 定义的收敛方式在当时是完全新的。
** 必须指出，直到 1935 年前后，「连续」函数一般指的是把每个收敛序列变为收敛序列的映射。
*** 对 Hilbert 而言，一个双线性型 \(B(x, y)\) 是完全连续的，如果只要序列 \((x_{n})\)、\((y_{n})\) 分别弱收敛于 \(x\) 与 \(y\)，\(B(x_{n}, y_{n})\) 就收敛于 \(B(x, y)\)。

与此同时，在 20 世纪的最初几年里，另外一些思想潮流出现并加强了导向赋范空间理论的趋势。「泛函」的一般概念（即定义在一个集上的数值函数，该集的元素本身是一个或多个实变量的数值函数）在 19 世纪的最后几十年里，一方面与变分法、另一方面与积分方程论相联系而得到复苏。但「算子」的一般概念主要是从以 Pincherle 为核心、尤其是以 Volterra 为核心的意大利学派产生的。这个学派的工作常常停留在相当形式的层面，且与特殊问题相关联，因为缺乏对底层拓扑概念足够深入的分析。1903 年，Hadamard 开创了「拓扑」对偶的现代理论，他在紧区间上连续数值函数的空间 \(\mathcal{C}(\mathrm{I})\)（赋以一致收敛拓扑）上寻求最一般的连续线性「泛函」，并把它们刻画为积分序列 \(x\mapsto \int_{\mathrm{I}}k_n(t)x(t)dt\) 的极限。

1907 年，Fréchet 与 F. Riesz 类似地证明了 Hilbert 空间上的连续线性型就是 Hilbert 引入的「有界」线性型；然后在 1909 年，F. Riesz 把 \(\mathcal{C}(\mathrm{I})\) 上的每个连续线性型表为一个 Stieltjes 积分，从而把 Hadamard 定理置于最终的形式，这一定理很久以后成了现代积分论的出发点（见 INT, chap.~II-V 的历史注记）。

第二年，F. Riesz (IX a) 又通过引入并研究空间 \(L^{p}(I)\)（仿照 Hilbert 空间的理论建立；它是区间 I 上 p 次幂可积的函数的空间，指数 p 满足 \(1 < p < +\infty\)）而在这一理论中取得了新的重要进展；三年后，这一研究又由关于序列空间 \(\ell^{p}(\mathbf{N})\) 的类似工作 (IX c) 所延续。正如我们后将看到的，这些研究对对偶思想的分类做出了很大贡献，因为我们第一次遇到了两个并非自然同构的对偶空间 *。

此后，F. Riesz 考虑过一个能涵盖所有这些结果的公理化研究 ((IX a), p.~452)，而且看来只是一位不愿偏离经典数学的分析学家的顾虑，才使他没有把他 1918 年那篇关于 Fredholm 理论的著名论文 (IX d) 写成这种形式。在那篇论文中他主要考虑紧区间上连续函数的空间 \(\mathcal{C}(\mathrm{I})\)；但在定义了这个空间的范数并注意到赋以该范数的 \(\mathcal{C}(\mathrm{I})\) 是完备的之后，他在论证中所用的就只有完备赋范空间的公理 **。我们不打算详细审查这一工作，仅指出：完全连续映射的概念（以把一个邻域变为相对紧集为其性质）在这项工作中第一次以一般的方式被定义 ***；通过公理化分析的一个杰作，Fredholm 的整个理论（就其定性方面而言）被化为单独一条基本定理，即每个局部紧的赋范空间都是有限维的。
* 尽管 \(L^1\) 与 \(L^\infty\) 之间的对偶已隐含在这个时期关于 Lebesgue 积分的大多数著作中，但直到 1918 年，H. Steinhaus 才证明了 \(L^1(\mathrm{I})\)（I 为有穷区间）上的每个连续线性型都具有形式 \(x\mapsto\int_{\mathrm{I}} f(t)x(t)dt\)，其中 \(f\in L^\infty(\mathrm{I})\)。
** 不过 F. Riesz 明确指出，他的定理对连续函数的应用只是远为一般的概念的一块「试金石」((IX d), p.~71)。
*** 在他关于 \(L^p\) 空间的工作中，F. Riesz 把完全连续映射定义为把每个弱收敛序列变为强收敛序列的映射；这（由于 \(L^p\)（对 \(1 < p < +\infty\)）中单位球的弱紧性）在这一情形等价于上述定义；此外，F. Riesz 指出，对 \(L^2\) 空间，他的定义等价于 Hilbert 的定义（把线性映射的语言翻译为双线性型的语言即可，(IX a), p.~487）。

赋范空间的一般定义是 1920 至 1922 年间由 S. Banach、H. Hahn 与 E. Helly 给出的（后者只考虑实数或复数的序列空间）。在随后的十年里，这些空间的理论主要围绕两个对应用具有根本重要性的问题发展：对偶论以及与 Baire「纲」概念相联系的定理。

我们已经看到，（拓扑意义上的）对偶思想起源于 20 世纪初；它是 Hilbert 理论的底层概念，并在 F. Riesz 的工作中占据中心地位。例如，后者在 1911 年 ((IX b), p.~41-42) 注意到，关系式 \(|f(x)| \leqslant M \| x\|\)（在 Hilbert 空间中取作「有界」线性泛函的定义）等价于 \(f\) 的连续性（在空间 \(\mathcal{C}(\mathrm{I})\) 的情形），而且这是用相当一般的论证证明的。关于 \(\mathcal{C}(\mathrm{I})\) 上连续线性泛函的刻画，他进一步注意到，一个集 A 在 \(\mathcal{C}(\mathrm{I})\) 中为完全集的条件是：在 I 上不存在与 A 中所有函数都「正交」的 Stieltjes 测度 \(\mu \neq 0\)（从而推广了 Gram 对完全标准正交系的条件）；最后，在同一工作中，他确立了空间 \(L^{\infty}\) 的对偶「大于」Stieltjes 测度组成的空间 ((IX b), p.~62)。

另一方面，F. Riesz 在他关于空间 \(L^{p}(\mathbf{I})\) 与 \(\ell^{p}(\mathbf{N})\) 的工作中，成功地修改了 E. Schmidt (VIII b) 给出的 Hilbert 空间中线性方程组的解法，使之适用于更一般的情形。E. Schmidt 的想法在于：在由方程组 (6) 所表示的闭线性流形中找出一个到原点的距离为极小的点，以此来确定 (6) 的一个「极值」解。利用同一想法，F. Riesz 证明了：存在一个函数 \(x \in \mathrm{L}^{p}(a, b)\) 满足方程

\[
\int_ {a} ^ {b} \alpha_ {i} (t) x (t) d t = b _ {i} (i = 1, 2, \dots)\tag{9}
\]

(其中 \(\alpha_{i}\) 属于 \(\mathbf{L}^q\)（满足 \(\frac{1}{p} +\frac{1}{q} = 1\)），且此外还满足 \(\int_{a}^{b}|x(t)|^{p}dt\leqslant \mathbf{M}^{p}\))的必要充分条件是：对实数的每个有穷序列 \((\lambda_i)_{1\leqslant i\leqslant n}\)，我们有

\[
\left| \sum_ {i = 1} ^ {n} \lambda_ {i} b _ {i} \right| \leqslant \mathrm{M} \left(\int_ {a} ^ {b} \left| \sum_ {i = 1} ^ {n} \lambda_ {i} \alpha_ {i} (t) \right| ^ {q} d t\right) ^ {1 / q}.\tag{10}
\]

1911 年 (IX b)，他以类似的方式处理了「广义矩量问题」，它在于求解方程组

\[
\int_ {a} ^ {b} \alpha_ {i} (t) d \xi (t) = b _ {i} (i = 1, 2, \dots)\tag{11}
\]

其中 \(\alpha_{i}\) 连续而未知量是一个 Stieltjes 测度 \(\xi\) *；在这一情形，问题显然可以重新表述为：由一个连续线性泛函在 \(\mathcal{C}(\mathrm{I})\) 中一个给定点序列上的值来确定这个泛函。Helly 在 1912 年就是以这种形式处理这一问题的 --- 以一种相当不同且适用范围广得多的方法得到了 F. Riesz 的条件 * --- 并于 1921 年以远为一般的条件再次处理了它。如上所见，他（在序列空间上）引入了范数的概念，并注意到这一概念推广了 n 维空间中凸体的「度规」，后者曾被 Minkowski 在他关于「数的几何」的著名著作 (IV) 中使用。在研究过程中，Minkowski 还（在 \(R^{n}\) 中）定义了支撑超平面与「支撑函数」的概念 (IV b)，并证明了在凸体边界的每一点处都存在支撑超平面 ((IV a), p.~33-35)。Helly 把这些概念推广到赋以任意范数的序列空间 E；他在 E 与空间 \(E'\)（由满足如下条件的序列 \(u = (u_{n})\) 组成）之间建立了一个对偶：对所有 \(x = (x_{n}) \in E\)，级数 \((u_{n}x_{n})\) 收敛；用 \(\langle u, x \rangle\) 表示这个级数的和，他在 \(E'\) 中由公式 \(\sup_{x \neq 0} |\langle u, x \rangle| / \|x\|\) 定义了一个范数，它在有限维空间中给出支撑函数 **。然后 Helly 证明了：E 中一个方程组 (6)（其中每个序列 \(u_{i} = (a_{ij})_{j \geqslant 1}\) 假定属于 \(E'\)）的求解可化为下列两个问题的依次求解：1. 在赋范空间 \(E'\) 上找一个连续线性型 L，使得对每个指标 i 有 \(\mathrm{L}(u_{i}) = b_{i}\)；如他所指出的，这导致 (10) 型的条件；2. 判明这样一个线性型能否写成 \(u \mapsto \langle u, x \rangle\)（对某个 \(x \in E\)）的形式。他注意到，后一问题即使 L 存在也未必有解，而且在某些特殊情形，他给出了几个蕴含解 \(x \in E\) 存在的充分条件 (X)。
* 经典的「矩量问题」对应于区间 \(]a, b[\) 为 \(]0, +\infty[\) 或 \(]-\infty, +\infty[\) 且 \(\alpha_{i}(t) = t^{i}\) 的情形；此外，还假定测度 \(\xi\) 是正的（F. Riesz 在他 1911 年的论文中指出，当寻求这种性质的解时，他的一般条件应如何修改）。在经典矩量问题的各种解法中，我们特别提到 F. Riesz 的方法，他非常优雅地把泛函演算的一般思想与单复变量函数论结合起来，得到了关于 \(b_{i}\) 的显式条件。(Sur le problème des moments, 3, Ark. för Math., t. XVII (1922-1923), no 16, 52 p.)
* 与 F. Riesz 一样 ((IX b), p.~49-50)，Helly 在他的证明中使用了一个「选择原理」，它恰好就是 Stieltjes 测度空间中单位球的弱紧性；F. Riesz 也曾在 \(L^p\) 空间 (1 < p < +\infty) 中使用过类似的性质。
** 为能以这种方式得到一个范数，必须假定关系式 \(\langle u,x\rangle = 0\)（对所有 \(x \in E\)）蕴含 \(u=0\)，Helly 明确指出了这一点。
*** Banach 早在 1923 年就给出过一个类似的论证，用以定义平面上的一个不变测度（对每个有界子集都有定义）(XII a)。

1927 年，这些思想在 H. Hahn 的一篇基础性论文 (XI) 中获得了它们的最终形式，其结果两年后由 S. Banach（独立地）重新发现 (XII b)。Hahn 把 Minkowski-Helly 的方法应用于任意的赋范空间，从而在对偶空间上定义了一个（完备）赋范空间的结构；这立即使得 Hahn 可以考虑一个赋范空间的逐次对偶，并以一般的方式提出自反空间的问题，正如 Helly 已经预见到的那样。但最重要的是，不增大范数地延拓一个连续线性泛函这一主要问题由 Hahn 一般地彻底解决，所用的论证是对维数做超限归纳 --- 由此给出了选择公理在泛函分析中的一个重要应用的最早例子之一 ***。在这些结果之外，Banach 又对连续线性映射与其转置之间的关系做了详细研究，借助一个关于对偶空间中弱闭子集的深刻定理（cf.~IV, p.~25, cor. 2），把先前只在 \(\mathbf{L}^p\) 空间情形已知的结果 (IX a) 推广到一般赋范空间；利用几年后由 Hausdorff 与 Banach 本人引入的赋范空间商空间的概念，这些结果可以表达得更为醒目。最后，发现单位球的弱紧性（如上所述曾在若干特殊情形中被观察到）与自反性之间关系的，再一次是 Banach，至少对满足第一可数公理的空间是如此。从那时起，赋范空间对偶论的主要轮廓就可以被认为已经确定。

在同一时期，一些表面上看似悖论的定理得到了阐明，它们的最初例子起源于 1910 年前后。那一年，Hellinger 与 Toeplitz 本质上证明了：Hilbert 空间上一个有界双线性型的序列 \(\mathbf{B}_{n}(x,y)\)，若其值 \(\mathbf{B}_{n}(a,b)\)（对每对 \((a,b)\)）都有界（界数先验地依赖于 a 与 b），则实际上在每个球上一致有界。他们的证明基于一个归谬论证，即归纳地构造出一个破坏假设的特殊对 \((a,b)\)；这一方法从此以「滑动峰」之名著称，并且在许多类似问题中仍然有用（cf.~IV, p.~54, exerc. 15）。1905 年，Lebesgue 曾用类似的方法证明存在 Fourier 级数在某些点发散的连续函数；而在与 Hellinger 和 Toeplitz 同年，他再次使用这一方法，证明了 \(L^{1}\) 中一个弱收敛序列在范数下有界 *。这些例子在随后几年里成倍增加，但直到 1927 年都没有引入任何新思想；那时 Banach 与 Steinhaus（在 Saks 的部分合作下）把这些现象与贫集概念以及完备度量空间中的 Baire 定理联系起来，从而得到了一个涵盖此前所有特殊情形的一般论断 (XIII)。同一时期，完备赋范空间中「纲」问题的研究使 Banach 得到了关于连续线性映射的若干其他结果；最引人注目也肯定最深刻的是「闭图像」定理，它与 Banach-Steinhaus 定理一样，已成为现代泛函分析的一个至关重要的工具 (XII b)。

Banach 关于「线性算子」的专著的出版 (XII c) 标志着赋范空间理论的成年。上面提到的所有结果以及许多其他结果都可以在这部书中找到，虽然组织得有些杂乱，但配有取自分析各个领域的大量引人注目的例子，这些例子似乎预示着这一理论的光明前景。这部著作取得了巨大的成功，其直接效果之一就是 Banach 所用的语言与记号几乎被普遍采用。但尽管过去 40 年间对 Banach 空间 (XVII) 进行了大量研究，如果排除 Banach 代数论及其对交换与非交换调和分析的应用，那么这一理论对经典分析重大问题几乎完全缺乏新的应用，这多少削弱了基于它的希望。

最富成果的发展更多是在拓宽与同赋范空间概念相关的更深入的公理化分析的方向上发生的。
* 还要注意 Landau 于 1907 年证明的一个类似（且更容易）的定理，它曾是 F. Riesz 的 \(L^p\) 空间理论的出发点：如果通项为 \(u_n x_n\) 的级数对每个序列 \((x_n)\in \ell^p(\mathbf N)\) 都收敛，则序列 \((u_n)\) 属于 \(\ell^q(\mathbf N)\)（其中 \(\frac{1}{p}+\frac{1}{q}=1\)）。







尽管自 20 世纪初以来所遇到的函数空间一般说来都带有「自然」范数，也存在某些例外。1910 年前后，E. H. Moore 提出了一致收敛概念的一个推广，用「相对一致收敛」的概念来代替它，其中 0 的一个邻域由函数 \(f\)（满足关系 \(|f(t)| \leqslant \varepsilon g(t)\)）组成，\(g\) 是一个处处 \(>0\) 且可随邻域而变的函数。另一方面，1930 年以前，人们已经注意到诸如简单收敛、可测函数的依测度收敛或整函数的紧收敛这样的概念不能借助范数来定义；1926 年，Fréchet 注意到这类向量空间可以是可度量化且完备的。但这些更一般空间的理论只有与凸性思想相联系才能富有成果地发展。后者（已出现在 Helly 的工作中）是 Banach 及其学生们所做研究的对象，他们认识到把赋范空间论的若干结果几何地解释的可能性，从而为 Kolmogoroff 与 J. von Neumann 在 1935 年给出的局部凸空间的一般定义铺平了道路。这些空间的理论、特别是与对偶相关的问题，主要是在 1950 年代发展起来的，我们在本册书中给出了这一研究的本质性结果。在这方面我们必须指出，一方面是一般拓扑基本概念所带来的简洁性与一般性方面的进步，这些概念是在 1930 至 1940 年间发展起来的；另一方面是有界集概念的重要性，它由 Kolmogoroff 与 von Neumann 于 1935 年引入，其在空间对偶理论中的根本作用由 Mackey (XIV) 与 Grothendieck (XVIII) 的工作揭示出来。最后，可以肯定，推动这些研究的主要动力来自向 Banach 理论不适用的那些领域的分析学应用的新可能性：在这方面，我们提到自 1934 年起由 Köthe、Toeplitz 及其学生们在一系列论文中发展的序列空间理论 (XV)，对 Fantappié 的「解析泛函」论的关注，尤其是 L. Schwartz 的分布理论 (XVI)，局部凸空间的现代理论在其中找到了一个远未枯竭的应用领域。

\backmatterentry{参考文献}

\begin{enumerate}
\def\labelenumi{(\Roman{enumi})}
\item
  C. STURM : a) Sur les équations différentielles linéaires du second ordre, Journ. de Math. (1), t. I (1836), p.~106-186 ; b) Sur une classe d'opérations à différences partielles, ibid., p.~373-444.
\item
  J. LIOUVILLE : a) Sur le développement des fonctions ou parties de fonctions en séries dont les divers termes sont assujettis à satisfaire à une même équation différentielle du second ordre contenant un paramètre variable, Journ. de Math. (1), t. I (1836), p.~253-265, t. II (1837), p.~16-35 et 418-436 ; b) D'un théorème dû à M. Sturm et relatif à une classe de fonctions transcendants, ibid., t. I (1836), p.~269-277.
\item
  J. P. GRAM, Ueber die Entwicklung reeller Functionen in Reihen mittelst der Methode der kleinsten Quadrate, J. de Crelle, t. XCIV (1883), p.~41-73.
\item
  H. MINKOWSKI : a) Geometrie der Zahlen, 1 \(^{re}\) éd., Leipzig (Teubner), 1896 ; b) Theorie der konvexen Körper, Gesammelte Abhandlungen, t. II, p.~131-229, Leipzig-Berlin (Teubner), 1911. (Réimpression, New York (Chelsea), 1967.)
\item
  H. POINCARÉ : a) Sur les équations de la Physique mathématique, Rend. Palermo, t. VIII (1894), p.~57-156 (= Œuvres, t. IX, p.~123-196, Paris (Gauthier-Villars), 1954); b) La méthode de Neumann et le problème de Dirichlet, Acta Mathematica, t. XX (1896), p.~59-142 (= Œuvres, t. IX, p.~202-272, Paris (Gauthier-Villars), 1954).
\item
  I. FREDHOLM, Sur une classe d'équations fonctionnelles, Acta Mathematica, t. XXVII (1903), p.~365-390.
\item
  D. HILBERT, Grundzüge einer allgemeinen Theorie der linearen Integralgleichungen, New York (Chelsea), 1953 (= Gött. Nachr., 1904, 1905, 1906, 1910).
\item
  E. SCHMIDT : a) Zur Theorie der linearen und nichtlinearen Integralgleichungen. I. Teil : Entwicklung willkürlicher Funktionen nach Systemen vorgeschriebener, Math. Ann., t. LXIII (1907), p.~433-476 ; b) Ueber die Auflösung linearer Gleichungen mit unendlich vielen Unbekannten, Rend. Palermo, t. XXV (1908), p.~53-77.
\item
  F. RIESZ : a) Untersuchungen über Systeme integrierbarer Funktionen, Math. Ann., t. LXIX (1910), p.~449-497 ; b) Sur certains systèmes singuliers d'équations intégrales, Ann. Ec. Norm. Sup. (3), t. XXVIII (1911), p.~33-62 ; c) Les systèmes d'équations linéaires à une infinité d'inconnues, Paris (Gauthier-Villars), 1913 ; d) Ueber lineare Funktionalgleichungen, Acta Mathematica, t. XLI (1918), p.~71-98 ; e) Zur Theorie des Hilbertschen Raumes, Acta litt. ac scient. (Szeged), t. VII (1934-35), p.~34-38.
\item
  E. HELLY, Ueber Systeme linearer Gleichungen mit unendlich vielen Unbekannten, Monatshefte für Math. und Phys., t. XXXI (1921), p.~60-91.
\item
  H. HAHN, Ueber lineare Gleichungssysteme in linearen Räumen, J. de Crelle, t. CLVII (1927), p.~214-229.
\item
  S. BANACH et H. STEINHAUS, Sur le principe de condensation des singularités, Fund. Math., t. IX (1927), p.~50-61.
\item
  S. BANACH : a) Sur le problème de la mesure, Fund. Math., t. IV (1923), p.~7-33 ; b) Sur les fonctionnelles linéaires, Studia Math., t. I (1929), p.~211-216 et 223-239 ; c) Théorie des opérations linéaires, Warszawa, 1932. (Réimpression, New York (Chelsea), 1963.)
\item
  G. W. MACKEY : a) On infinite-dimensional linear spaces, Trans. Amer. Math. Soc., t. LVII (1945), p.~155-207 ; b) On convex topological spaces, Trans. Amer. Math. Soc., t. LX (1946), p.~519-537.
\item
  G. KÖTHE, Neubegründung der Theorie der vollkommenen Räume, Math. Nachr., t. IV (1951), p.~70-80.
\item
  L. SCHWARTZ, Théorie des distributions, 2 \(^{e}\) édition, Paris (Hermann), 1966.
\item
  J. LINDENSTRAUSS and L. TZAFRIRI, Classical Banach spaces, t. I, Berlin-Heidelberg-New York (Springer), 1977.
\item
  A. GROTHENDIECK : a) Produits tensoriels topologiques et espaces nucléaires, Mem. Amer. Math. Soc., n° 16 (1955); b) Espaces vectoriels topologiques, 3e éd., São Paulo (Publ. Soc. Mat. São Paulo), 1964.
\end{enumerate}

\backmatterentry{记号索引}

引用编号指明章与页（偶尔也指习题）

\(|\xi|\) , \(\|x\|\) : I, p.~3.

\(\mathcal{B}(I; K)\) , \(\mathcal{B}_{K}(I)\) , \(\ell_{K}^{\infty}(I)\) , \(\ell_{K}^{1}(I)\) , \(\mathcal{B}(I)\) , \(\ell^{1}(I)\) : I, p.~4.

\(E_{A}\)（A 为实向量空间 E 中的对称凸集）：II, p.~26.

\(\langle x, y \rangle\) : II, p.~42.

\(\sigma(F, G)\) : II, p.~42.

\(M^{\circ}\) , \(M^{\circ\circ}\) : II, p.~44.

\(^{t}u\)（u 为线性映射）：II, p.~46.

\(\mathcal{R}(X)\) : III, p.~9.

\(\mathcal{C}^{\infty}(U)\) : III, p.~9.

\(\mathcal{C}_{H}^{\infty}(U)\) , \(\mathcal{C}_{0}^{\infty}(U)\) : III, p.~9.

\(\mathcal{G}_{s,M}(I)\) , \(\mathcal{G}_{s}(I)\) , \(\mathcal{C}(I)\) : III, p.~10.

\(\mathcal{H}(U)\) , \(\mathcal{H}(L)\)（U 为 \(C^{n}\) 的开子集，L 为 \(C^{n}\) 的紧子集）：III, p.~10.

\(\mathcal{L}(E; F)\) : III, p.~13.

\(\mathcal{L}_{\Xi}(E; F)\) : III, p.~13.

\(\mathcal{L}_{s}(E; F)\) , \(\mathcal{L}_{c}(E; F)\) , \(\mathcal{L}_{pc}(E; F)\) , \(\mathcal{L}_{cc}(E; F)\) , \(\mathcal{L}_{b}(E; F)\) : III, p.~14.

\(E'\) , \(E'_{\Xi}\) , \(E'_{s}\) , \(E'_{c}\) , \(E'_{pc}\) , \(E'_{cc}\) , \(E'_{b}\) : III, p.~14.

\(\mathcal{L}(E)\) , \(\mathcal{L}_{\Xi}(E)\) , \(\mathcal{L}_{s}(E)\) , \(\mathcal{L}_{c}(E)\) , \(\mathcal{L}_{pc}(E)\) , \(\mathcal{L}_{cc}(E)\) , \(\mathcal{L}_{b}(E)\) : III, p.~14.

\(p_{M}\)（p 为半范数，M 为有界子集）：III, p.~14.

\(\mathcal{C}_{0}(\mathbf{R})\) : III, p.~18.

\(\tau(E, F)\) : IV, p.~2.

\(\beta(E, F)\) : IV, p.~4.

\(c_{E}\) : IV, p.~14.

\(\ell_{0}^{\infty}(\mathbf{N})\) : IV, p.~17.

\(c_{0}(\mathbf{N})\) , \(\ell^{1}(\mathbf{N})\) : IV, p.~18.

S(E): IV, p.~26.

\(H_{p}\) : IV, p.~26.

\(E_{\sigma}\) : IV, p.~32.

\(\mathcal{C}_{s}(\mathbf{X})\) : IV, p.~33.

\(C^{b}(\mathbf{X})\) , \(C(\mathbf{X})\) : IV, p.~36.

\(\mathcal{B}(\mathbf{X};\mathbf{R})\) : IV, p.~40.

Ind(u)（u 为 Fredholm 算子）：IV, p.~66, exerc. 21.

\(\overline{\xi}\) : V, p.~1.

\(\ell^{2}\) , \(\ell^{2}(\mathbf{N})\) : V, p.~4.

\(E_{(\mathbf{C})}\) : V, p.~4.

\(\langle x|y\rangle\), \(\|x\| = \langle x|x\rangle^{1/2}\), \((x|y) = \langle y|x\rangle\): V, p.~5.

\(\overline{\mathrm{E}}\)（E 为复准希尔伯特空间）：V, p.~6.

\(\mathcal{H}^{s} (\text {Sobolev space}): V, p.~6.\) \(H^{2}(\text {D}): V, p.~7.\) \(\mathcal{C}_{0}^{1}(U): V, p.~8.\) \(p_{H} (\text {H convex separated and complete set in a prehilbertian space}): V, p.~10.\) \(x^{*} (x a vector of a Hilbert space): V, p.~15 et p.~40.\) \(\bigoplus_{i\in I}\mathrm{E}_{i},\bigoplus_{i\in I}\mathrm{E}_{i}:V,p.18.\) \(E_{1}\oplus E_{2}\oplus\cdots\oplus E_{n} (E_{i}\text { Hilbert spaces}):V,p.18.\) \(\ell_{E}^{2}(I),\ell^{2}(I)(E\text {a Hilbert space}):V,p.18.\)

\(E_{1} \otimes_{2} E_{2}, \|z\|_{2} (z \in E_{1} \otimes_{2} E_{2}) : V, p.~26.\)

\(E_{1} \otimes_{2} E_{2} \otimes_{2} \cdots \otimes_{2} E_{n}, \bigotimes_{i=1}^{n} E_{i}, \|z\|_{2} \left( z \in \bigotimes_{i=1}^{n} E_{i} \right) : V, p. 27.\)

\(E_{1} \hat{\otimes}_{2} E_{2} \hat{\otimes}_{2} \ldots \hat{\otimes}_{2} E_{n}, \bigotimes_{1 \leqslant i \leqslant n} \hat{\otimes}_{2} E_{i}: V, p. 28\)

\(u_{1}\hat{\otimes}_{2}u_{2}\hat{\otimes}_{2}\ldots \hat{\otimes}_{2}u_{n}(u_{i}\) 为线性映射）：V, p.~28.

\(\hat{\mathbf{T}}^n (\mathrm{E}),\mathrm{E}\hat{\otimes}_n:\mathrm{V},\) p.~29.

\(\hat{\mathbf{S}}^{n}(\mathrm{E}), \hat{\mathbf{S}}(\mathrm{E}) : \mathrm{V}, \mathrm{p.} 30.\)

\(\hat{\mathbf{A}}^{n}(\mathrm{E}), \hat{\mathbf{A}}(\mathrm{E}) : \mathrm{V}, \mathrm{p. 33}\) .

\(\hat{\mathbf{T}}^n (u),\hat{\mathbf{S}}^n (u)\) \((u\) 为线性映射）：V, p.~31 与 p.~32.

\(\hat{\Lambda}^n (u)\)（\(u\) 为线性映射）：V, p.~34.

v.u, vu（u、v 为线性映射）：V, p.~37.

\(u^{*}\)（u 为线性映射）：V, p.~38.

\(\mathcal{H}(\mathrm{E})\)（E 为希尔伯特空间）：V, p.~44.

\(u \geqslant 0\)（u 为希尔伯特空间的自同态）：V, p.~45.

\(\mathcal{L}_{+}(\mathrm{E}):\mathrm{V},\mathrm{p.}45.\)

\(u \geqslant v\)（u、\(v\) 属于 \(\mathcal{L}(\mathrm{E})\)，E 为希尔伯特空间）：V, p.~45.

\(\tau (u)\)（\(u\) 为有限秩自同态）：V, p.~48.

\(\operatorname{Tr}(u) (u \geqslant 0 \text{ in } \mathcal{L}(\mathrm{E})): \mathrm{V}, \text{ p. } 49.\)

\(\mathcal{L}^1 (\mathrm{E})\)（E 为希尔伯特空间）：V, p.~51.

\(\mathcal{L}^2 (\mathrm{E};\mathrm{F}),\mathcal{L}^2 (\mathrm{E})\)（E、F 为希尔伯特空间）：V, p.~52.

\(\| u\| _2\) \((u\in \mathcal{L}(\mathrm{E};\mathrm{F}),\mathrm{E},\mathrm{F}\) 为希尔伯特空间）：V, p.~52.

\(\mathrm{Tr}(\mathrm{Q} / \mathrm{H})\)（Q、H 为正二次型）：V, p.~57.

\backmatterentry{术语索引}

吸收集、一集对另一集的吸收：I, p.~7. 适配有界结构：III, p.~3. 伴随：V, p.~38. 仿射变换：IV, p.~39. 与拓扑向量空间相伴的 Hausdorff 向量空间：I, p.~4.

集的闭均衡凸包：II, p.~13 与 p.~62. 集的均衡核：I, p.~7. 集的均衡包：I, p.~7. 均衡集：I, p.~6. Banach-Dieudonné 定理：IV, p.~24. Banach-Saks-Kakutani 定理：V, p.~68, exerc. 33. Banach 空间：I, p.~5. Banach-Steinhaus 定理：III, p.~25. Banach 定理：I, p.~17. 桶：III, p.~24. 桶型空间：III, p.~24. 有界结构的基：III, p.~1. 希尔伯特空间的（代数）基：V, p.~22. （Banach）基：IV, p.~69, exerc. 14. （完全 Banach、收缩 Banach）基：IV, p.~70, exerc. 15. （标准正交）基：V, p.~22. （无条件 Banach）基：IV, p.~71, exerc. 16. Bessel 不等式：V, p.~21. 双对偶：IV, p.~14. 双极定理：II, p.~44. Birkhoff-Alaoglu 定理：V, p.~78, exerc. 13. Bishop-Phelps 定理：II, p.~77, exerc. 4. 有界型局部凸空间：III, p.~12. 有界结构：III, p.~1. （适配的）有界结构：III, p.~3. （典范的）有界结构：III, p.~3. （凸的）有界结构：III, p.~2. 由一集族生成的有界结构：III, p.~1. （乘积）有界结构：III, p.~2. 吸有界集的集：III, p.~39, exerc. 11. 有界集：III, p.~2 与 p.~37, exerc. 1.

典范有界结构：III, p.~3. 从 \(\bigoplus_{i\in I}E_i'\) 到 \(\left(\prod_{i\in I}E_i\right)'\) 中的典范映射：IV, p.~13. E 到 E'\,' 的典范映射：IV, p.~14. E 到 E' 上的典范映射（E 为希尔伯特空间）：V, p.~15. 有限维向量空间上的典范拓扑：I, p.~2. 帽：II, p.~57. Cauchy-Schwarz 不等式：V, p.~3. 闭图像定理：I, p.~19. 由闭超平面定义的闭半空间：II, p.~15. 上边缘：IV, p.~72, exerc. 3. 1-上闭链（连续的）：IV, p.~72, exerc. 3.

紧线性映射：III, p.~6. 序向量结构的相容拓扑：II, p.~15. 向量空间结构与预序的相容：II, p.~12. 向量空间结构与拓扑的相容：III, p.~1. 与对偶相容的（局部凸拓扑）：IV, p.~1. （正交）补：V, p.~13. 完备拓扑向量空间：I, p.~5. Hausdorff 准希尔伯特空间的完备化：V, p.~8. Hausdorff 拓扑向量空间的完备化：I, p.~6. 复线性型：II, p.~61. 复线性流形：II, p.~61. 复局部凸空间：II, p.~62. 复化（准希尔伯特空间的）：V, p.~5. 复化的拓扑向量空间：II, p.~62. 凹函数：II, p.~17. （渐近）锥：II, p.~67. 由一集生成的（凸）锥：II, p.~11. （尖的与非尖的）锥：II, p.~10. （多面）锥：II, p.~91. （真尖凸）锥：II, p.~11. 复准希尔伯特空间的共轭：V, p.~6. 集的均衡凸包：II, p.~10 与 p.~62. 凸有界结构：III, p.~2. 集的闭凸包：II, p.~13. 凸函数：II, p.~17. 凸集：II, p.~7 与 p.~62. 集的（对称）凸包：II, p.~16. 关于标准正交基的坐标：V, p.~22. 集的（均衡）核：I, p.~7.

序的密度：V, p.~7. （希尔伯特）维数：V, p.~24. 凸集的维数：II, p.~10. Dirichlet 空间：V, p.~8. 远距集：IV, p.~72, exerc. 1. 区分空间：IV, p.~52, exerc. 4. 实拓扑向量空间的（代数）对偶：II, p.~42. 局部凸空间（实的或复的）的对偶：III, p.~14. 实拓扑向量空间的对偶：II, p.~42. （弱、强）对偶：III, p.~14. 在 F 中分离的对偶、分离对偶：II, p.~41. （向量空间间的）对偶：II, p.~40. Dvoretzky-Rogers 定理：V, p.~63, exerc. 14.

Eberlein 定理：IV, p.~35. D. Edwards 定理：II, p.~94, exerc. 41. 自同态（自伴的）：V, p.~44. 自同态（正规的）：V, p.~43. 自同态（正的）：V, p.~45. 集的（均衡）包：I, p.~7. 集的（均衡凸）包：II, p.~10. 集的（均衡凸闭）包：II, p.~13 与 p.~62. 集的（闭凸）包：II, p.~13. 集的（凸）包：II, p.~9. 集的（对称凸）包：II, p.~10. 集的（对称凸闭）包：II, p.~13. S-等度亚连续、T-等度亚连续、(S, T)-等度亚连续集：III, p.~47, exerc. 7.

滤过化半范数集：II, p.~3.

（拓扑无关）族：I, p.~11.

Hausdorff 局部凸空间的穷竭、可穷竭空间：III, p.~49, exerc. 1.

凸锥的极生成元：II, p.~57.

凸集的极点：II, p.~54.

刻面：II, p.~87, exerc. 3.

（对偶）刻面：II, p.~87, exerc. 6.

（标准正交）族：V, p.~21.

终局部凸拓扑：II, p.~27.

终子空间：V, p.~41.

Fock 空间：V, p.~32 与 p.~34.

使两空间成对偶的（双线性）型：II, p.~40.

（复线性、实线性）型：II, p.~61.

（埃尔米特）型：V, p.~1.

（正埃尔米特）型：V, p.~2.

（正线性）型：II, p.~13.

与埃尔米特型相伴的（非退化埃尔米特）型：V, p.~2.

Fréchet 空间：II, p.~24 与 p.~63.

（凹、凸、严格凹、严格凸）函数：II, p.~16-17.

（正定）函数：V, p.~8.

（正齐次）函数、次线性函数：II, p.~19-20.

基本半范数系：II, p.~3.

凸集的度规：II, p.~20.

高斯空间：V, p.~32.

由一集族生成的（有界结构）：III, p.~1.

凸锥的（极）生成元：II, p.~57.

Gevrey 空间：III, p.~10.

Gram 行列式：V, p.~71, exerc. 7.

Grothendieck 定理：III, p.~20.

可赋均值的群：IV, p.~72, exerc. 4.

Haar 定理：II, p.~83, exerc. 8.

Hadamard 不等式：V, p.~37.

Hahn-Banach 定理：II, p.~22, p.~36 与 p.~63.

由闭超平面确定的（闭、开）半空间：II, p.~15.

Hardy 空间：V, p.~7.

拓扑向量空间的 Hausdorff 完备化：I, p.~6.

与拓扑向量空间相伴的 Hausdorff 向量空间：I, p.~4.

Helly 定理：II, p.~68, exerc. 21.

自伴算子：V, p.~44.

Hilbert-Schmidt 映射：V, p.~52.

Hilbert 空间、希尔伯特空间：V, p.~6.

集的（支撑）超平面：II, p.~37.

S-亚连续、T-亚连续、(S, T)-亚连续双线性映射：III, p.~30.

Fredholm 算子的指标：IV, p.~66, exerc. 21.

向量子空间上诱导的准希尔伯特结构：V, p.~6.

局部凸空间或拓扑的归纳极限：II, p.~29.

（Hadamard）不等式：V, p.~37.

（Bessel）不等式：V, p.~21.

（Cauchy-Schwarz）不等式：V, p.~3.

次桶型空间：III, p.~44, exerc. 7.

初始子空间：V, p.~41.

初始拓扑：IV, p.~4.

凸集的代数内点：II, p.~26.

James-Klee 定理：IV, p.~57, exerc. 25. R. C. James 空间：IV, p.~71, exerc. 18.

Krein-Milman 定理：II, p.~55. Krein 定理：IV, p.~37.

Legendre 多项式：V, p.~24. 复局部凸空间：II, p.~62. 实局部凸空间：II, p.~23. 局部凸拓扑：II, p.~23 与 p.~62.

Mackey 定理：IV, p.~2. Mackey 拓扑：IV, p.~2. （Hilbert-Schmidt）映射：V, p.~52. Markoff-Kakutani 定理：IV, p.~39. （Hilbert）矩阵：V, p.~75, exerc. 3. 关于标准正交基的矩阵：V, p.~22. 均值：IV, p.~40. 可度量化拓扑向量空间：I, p.~16. 极小型空间：II, p.~85, exerc. 13. Montel 空间：IV, p.~18.

正规算子：V, p.~43.

由闭超平面定义的开半空间：II, p.~15. （Fredholm）算子：IV, p.~66, exerc. 21. （酉）算子：V, p.~41. 序向量空间：II, p.~12. 对偶空间中子空间的正交：II, p.~44. 正交投影算子：V, p.~13. 准希尔伯特空间中的正交集：V, p.~13. 对偶空间中的正交集：II, p.~41. 准希尔伯特空间中与一子集正交者：V, p.~13. 一个对偶中的正交向量：II, p.~41. 准希尔伯特空间中的正交向量：V, p.~5. 标准正交基：V, p.~22. 标准正交集、族：V, p.~21. 标准正交化：V, p.~24. （初始、终）正交投影：V, p.~41.

抛物凸集：II, p.~67, exerc. 17. Parseval 等式：V, p.~22. 部分等距映射：V, p.~42. 尖锥：II, p.~10. 点（暴露的、严格凸的）：II, p.~88, exerc. 6. 超平面同侧的点、严格同侧的点：II, p.~9. Hilbert-Schmidt 自同态的极分解：V, p.~79, exerc. 14. 集的极集：II, p.~44 与 64. 极化公式：V, p.~2. 多面体：II, p.~90, exerc. 24. 正算子：V, p.~45. 正埃尔米特型：V, p.~2. 正齐次函数：II, p.~19. 预对偶：IV, p.~56, exerc. 23. 准希尔伯特半范数：V, p.~4. 准希尔伯特空间：V, p.~4. 预序向量空间：II, p.~12. 奇点凝聚原理：III, p.~42, exerc. 10.

乘积有界结构：III, p.~2. 到凸集上的投影：V, p.~11. （正交）投影：V, p.~13. 真尖凸锥：II, p.~11. 勾股定理：V, p.~12.

拟完备空间：III, p.~8.

实线性型、实线性流形：II, p.~61. 实局部凸空间：II, p.~23. 自反空间：IV, p.~16. 相对有界空间：III, p.~43, exerc. 6. （酉）表示：IV, p.~44. Ryll-Nardzewski 定理：IV, p.~43.

内积：V, p.~5. 标量平方：V, p.~5. 线段（闭的、开的、在 x 处开且在 y 处闭的）：II, p.~7. 准希尔伯特空间的半自同构：V, p.~6. 半桶型空间：IV, p.~21. 半完备空间：III, p.~7. 半范数：II, p.~1. 半赋范空间：II, p.~2. （准希尔伯特）半范数：V, p.~4. 半自反空间：IV, p.~15. 准希尔伯特空间的分离完备化：V, p.~8. 拓扑向量空间的分离完备化：I, p.~6. 被（闭超平面）分离的（集）：II, p.~37. 分别连续双线性映射：III, p.~28. 分别等度连续：III, p.~47. 分离对偶：II, p.~41. 超平面的（同侧、严格同侧的）点：II, p.~8. 单纯形：II, p.~71, exerc. 41. Šmulian 定理：IV, p.~36. Sobolev 空间：V, p.~6. 锥的底：II, p.~60. （DF）空间：IV, p.~57, exerc. 2. （拓扑向量）空间：I, p.~1. （弱）空间：II, p.~42. （标量）平方：V, p.~5. 星形集：II, p.~65, exerc. 1. 局部凸空间序列、拓扑序列的严格归纳极限：II, p.~33. 严格凹、严格凸函数：II, p.~16-17. 被（闭超平面）严格分离的（集）：II, p.~37. 强对偶：III, p.~14. E' 的强有界子集：III, p.~14. 次线性函数：II, p.~20. 连续线性映射的（终、初始）子空间：V, p.~41. （准希尔伯特）子空间：V, p.~6. 希尔伯特空间的（外希尔伯特）和：V, p.~18. 向量子空间的（希尔伯特）和：V, p.~18. 局部凸空间或拓扑的（拓扑直）和：II, p.~30. 支撑流形：II, p.~87, exerc. 3. 集的对称闭凸包：II, p.~13.

Tchebycheff 定理：II, p.~84, exerc. 8. （希尔伯特）张量积：V, p.~28. 准希尔伯特空间的张量积：V, p.~26-27.

拓扑无关集、族：I, p.~12 与 p.~11.

\(\mathfrak{S}\) -拓扑：III, p.~13 与 IV, p.~2.

与序向量空间结构相容的拓扑：II, p.~15.

与向量空间结构相容的拓扑：I, p.~1.

由半范数、由一半范数组定义的拓扑：II, p.~2-3.

（初始）拓扑：IV, p.~4.

（局部凸）拓扑：II, p.~23 与 p.~62.

（Mackey）拓扑：IV, p.~2.

简单收敛、紧收敛、预紧收敛、紧凸收敛、有界收敛拓扑：III, p.~14.

（弱）拓扑：II, p.~42.

（弱化）拓扑：IV, p.~4.

完全集：I, p.~11.

连续正自同态的迹：V, p.~49.

一个正二次型关于另一个的迹：V, p.~57.

（仿射）变换：IV, p.~39.

连续线性映射的转置：II, p.~46 与 IV, p.~6.

超有界型空间：III, p.~45, exerc. 19.

超范数：I, p.~26, exerc. 12.

超半范数：II, p.~2.

酉算子：V, p.~41.

酉表示：IV, p.~44.

无尖锥：II, p.~10.

（复线性、实线性）流形：II, p.~61.

（支撑）流形：II, p.~87, exerc. 3.

弱对偶：III, p.~14.

弱有界：III, p.~14.

弱拓扑、弱空间：II, p.~42.

弱化拓扑：IV, p.~4.

\backmatterentry{Banach 空间若干重要性质的综述}

为方便读者起见，我们在此汇集赋范空间、特别是 Banach 空间的主要结果。标量域 K 为 R 或 C。

线性映射空间；对偶

\begin{enumerate}
\def\labelenumi{\arabic{enumi})}
\tightlist
\item
  设 \(\mathbf{E}\) 与 \(\mathbf{F}\) 为两个赋范空间。线性映射 \(u\)（从 \(\mathbf{E}\) 到 \(\mathbf{F}\) 中的）为连续的，当且仅当
\end{enumerate}

\[
\| u \| = \sup _ {\| x \| \leqslant 1} \| u (x) \|\tag{1}
\]

为有限的。映射 \(u \mapsto \| u \|\) 是向量空间 \(\mathcal{L}(\mathrm{E};\mathrm{F})\)（由 \(\mathrm{E}\) 到 \(\mathrm{F}\) 中的连续线性映射所组成的）上的一个范数。

设 \(\mathbf{F}\) 为 Banach 空间。则 \(\mathcal{L}(\mathrm{E};\mathrm{F})\) 是 Banach 空间。完备化 \(\hat{\mathrm{E}}\)（\(\mathrm{E}\) 的）是 Banach 空间，且映射 \(u\mapsto u|\mathrm{E}\) 是从 \(\mathcal{L}(\hat{\mathrm{E}};\mathrm{F})\) 到 \(\mathcal{L}(\mathrm{E};\mathrm{F})\) 上的一个双射等距。

\begin{enumerate}
\def\labelenumi{\arabic{enumi})}
\setcounter{enumi}{1}
\tightlist
\item
  设 E 为赋范空间。记 \(E' = \mathcal{L}(E; K)\)，其中 K 赋有范数 \(\lambda \mapsto |\lambda|\)。Banach 空间 \(E'\) 称为 E 的对偶，而对偶 \(E''\)（\(E'\) 的）称为 E 的双对偶。
\end{enumerate}

用 \(\sigma(E, E')\) 表示 \(E\) 上使所有线性型 \(x' \in E'\) 都连续的最粗拓扑；它称为 \(E\) 的弱化拓扑。用 \(\sigma(E', E)\) 表示 \(E'\) 上使线性型 \(x' \mapsto \langle x', x \rangle\)（\(E'\) 上的、其中 \(x\) 取遍 \(E\) 的）都连续的最粗拓扑；这时 \(\sigma(E', E)\) 称为 \(E'\) 上的弱拓扑。由范数导出的 \(E'\) 上的拓扑称为强拓扑。

\begin{enumerate}
\def\labelenumi{\arabic{enumi})}
\setcounter{enumi}{2}
\tightlist
\item
  设 E 为赋范空间，M 为 E 的闭向量子空间。设 \(\pi\) 为从 E 到 E/M 上的典范映射。向量空间 E/M 上的一个范数由下式定义
\end{enumerate}

\[
\| \xi \| = \inf _ {\pi (x) = \xi} \| x \|.\tag{2}
\]

当 E 为 Banach 空间时，M 与 E/M 亦然。对每个赋范空间 F，线性映射 \(u \mapsto u \circ \pi\)（从 \(\mathcal{L}(\mathrm{E}/\mathrm{M}; \mathrm{F})\) 到 \(\mathcal{L}(\mathrm{E}; \mathrm{F})\) 中的）是等距的。

\begin{enumerate}
\def\labelenumi{\arabic{enumi})}
\setcounter{enumi}{3}
\tightlist
\item
  设 E 为赋范空间。对每个 \(x' \in E'\)，由定义我们有
\end{enumerate}

\[
\| x^{\prime}\| = \sup_{\substack{\| x\| \leqslant 1\\ x\in \mathbf{E}}}\left|\langle x^{\prime},x\rangle \right|.\tag{3}
\]

此外（Hahn-Banach 定理），我们有

\[
\| x \| = \sup_{\substack{\| x^{\prime}\| \leqslant 1\\ x^{\prime}\in \mathbf{E}^{\prime}}}\left|\langle x^{\prime},x\rangle \right|\tag{4}
\]

对所有 \(x \in E\) 成立。换言之，从 E 到其双对偶 \(E''\) 中的典范映射是等距的。
